\documentclass[11pt]{article}
\usepackage{amsmath,amssymb,array,booktabs,longtable,geometry,fontspec,xcolor,alltt}
\usepackage{tikz}\usetikzlibrary{arrows.meta}
\usepackage[hidelinks]{hyperref}
\definecolor{claimlink}{RGB}{25,70,150}
\newcommand{\pvar}[2]{\hyperlink{#1-P#2}{\textcolor{claimlink}{P_{#2}}}}
\newcommand{\pvartarget}[2]{\raisebox{\ht\strutbox}{\hypertarget{#1-P#2}{}}}
\makeatletter
\newenvironment{derivation}
  {\begin{list}{}{\setlength{\leftmargin}{2.6em}\setlength{\labelwidth}{2.2em}%
    \setlength{\labelsep}{0.4em}\setlength{\itemsep}{2pt}\setlength{\parsep}{0pt}%
    \setlength{\topsep}{6pt}\@beginparpenalty\@M\@itempenalty\@M%
    \interlinepenalty\@M\raggedright}}
  {\end{list}}
\newcommand{\breakrow}[1]{\@itempenalty\z@\item[#1]\@itempenalty\@M}
\newcommand{\keepwithnext}{\@endparpenalty\@M}
\makeatother
\clubpenalty=10000
\widowpenalty=10000
\newcommand{\gloss}[1]{\enspace---\enspace{\small\itshape #1}}
\geometry{margin=1in}
\setlength{\emergencystretch}{3em}
\frenchspacing
\setmainfont{FreeSerif}[Extension=.otf,UprightFont=*,BoldFont=*Bold,
  ItalicFont=*Italic,BoldItalicFont=*BoldItalic,Ligatures=TeX]
\setmonofont{DejaVuSansMono}[Extension=.ttf,UprightFont=*,BoldFont=*-Bold,
  ItalicFont=*-Oblique,BoldItalicFont=*-BoldOblique,Scale=MatchLowercase]
\newfontfamily\hebrewfont{FreeSerif}[Extension=.otf,UprightFont=*,BoldFont=*Bold,
  ItalicFont=*,BoldItalicFont=*Bold,Script=Hebrew]
\newfontfamily\syriacfont{FreeSans}[Extension=.otf,UprightFont=*,BoldFont=*Bold,
  ItalicFont=*,BoldItalicFont=*Bold,Script=Syriac]
\newfontfamily\symbolfont{FreeSerif}[Extension=.otf,UprightFont=*,BoldFont=*Bold,
  ItalicFont=*,BoldItalicFont=*Bold]
\XeTeXinterchartokenstate=1
\newcommand{\fontrange}[3]{\count255=#2\relax
  \loop\XeTeXcharclass\count255=#1\relax\ifnum\count255<#3\relax
  \advance\count255 by 1\repeat}
\newXeTeXintercharclass\hebrewclass
\fontrange\hebrewclass{"0590}{"05FF}
\newXeTeXintercharclass\syriacclass
\fontrange\syriacclass{"0700}{"074F}
\newXeTeXintercharclass\symbolclass
\fontrange\symbolclass{"1D2C}{"1D6A}
\fontrange\symbolclass{"22C0}{"22C3}
\XeTeXinterchartoks 0 \hebrewclass {\begingroup\hebrewfont}
\XeTeXinterchartoks 4095 \hebrewclass {\begingroup\hebrewfont}
\XeTeXinterchartoks \hebrewclass 0 {\endgroup}
\XeTeXinterchartoks \hebrewclass 4095 {\endgroup}
\XeTeXinterchartoks \hebrewclass \syriacclass {\endgroup\begingroup\syriacfont}
\XeTeXinterchartoks \hebrewclass \symbolclass {\endgroup\begingroup\symbolfont}
\XeTeXinterchartoks 0 \syriacclass {\begingroup\syriacfont}
\XeTeXinterchartoks 4095 \syriacclass {\begingroup\syriacfont}
\XeTeXinterchartoks \syriacclass 0 {\endgroup}
\XeTeXinterchartoks \syriacclass 4095 {\endgroup}
\XeTeXinterchartoks \syriacclass \hebrewclass {\endgroup\begingroup\hebrewfont}
\XeTeXinterchartoks \syriacclass \symbolclass {\endgroup\begingroup\symbolfont}
\XeTeXinterchartoks 0 \symbolclass {\begingroup\symbolfont}
\XeTeXinterchartoks 4095 \symbolclass {\begingroup\symbolfont}
\XeTeXinterchartoks \symbolclass 0 {\endgroup}
\XeTeXinterchartoks \symbolclass 4095 {\endgroup}
\XeTeXinterchartoks \symbolclass \hebrewclass {\endgroup\begingroup\hebrewfont}
\XeTeXinterchartoks \symbolclass \syriacclass {\endgroup\begingroup\syriacfont}
\usepackage{newunicodechar}
\newunicodechar{🡒}{→}
\title{The canonical witness — faith alone, from the texts all parties accept}
\author{\credit{TestimonyArguments.CanonicalWitness}}
\date{\editiondate\\[2pt]{\small Version \editionversion\ifeditionpreview\\\editionremark\fi}}
\newif\ifeditionrecorded
\newif\ifeditionpreview
\newcommand{\editiondate}{\today}
\newcommand{\editionversion}{unrecorded}
\newcommand{\editioncommit}{COMMIT}
\newcommand{\editionremark}{}
\InputIfFileExists{edition.def}{}{}
\newcommand{\credit}[1]{\ifcsname credit@#1\endcsname\csname credit@#1\endcsname\fi}
\newcommand{\creditline}[1]{\ifcsname credit@#1\endcsname{\centering\credit{#1}\par\medskip}\fi}
\InputIfFileExists{credits.def}{}{}
\begin{document}
\maketitle
\noindent This is one argument of the Testimony library, printed on its own. The library also prints every argument together, as one book (\texttt{arguments.pdf}). Each claim is given a numbered letter, $P_1$, $P_2$ and so on, so that the steps can be written out in the notation of logic, and a letter means the same claim wherever it appears in one argument. The table of numbered claims gives each in words, with who holds it and where they say so, and every letter links to its row.

\noindent Nothing here asserts that the premises are true. It shows only whether the conclusion follows from them. The last section says how this edition was made, and how to make it again.

\tableofcontents
\clearpage
\section*{The canonical witness — faith alone, from the texts all parties accept}
\addcontentsline{toc}{section}{The canonical witness — faith alone, from the texts all parties accept}

\subsection*{Faith alone from the texts every party accepts}
\addcontentsline{toc}{subsection}{Faith alone from the texts every party accepts}
The sola fide argument works through the disputes that surround the doctrine: what “works of the law” (ἔργα νόμου) denotes, whether “the faith of Christ” (πίστις Χριστοῦ) is our faith in him, what “saved” (σῴζω) means at Luke 7:50. This argument asks a narrower question, one those disputes can bury. Grant only the texts that Rome and the Reformers both hold as Scripture. What does each part of the New Testament say about how a person is justified, and does faith alone follow?

Every text is a premise cited to Scripture alone. Romans, Galatians, Ephesians and Titus are read together as Paul, whoever wrote the last two, because every party receives them as canon and the question is what the canon says. Hebrews is read as a corpus of its own, because it is anonymous. Trent grants all of them, and its Decree on Justification quotes most. So the texts are rated \texttt{consensus}, the highest rating, and nothing turns on them.

The argument is made in the readings. Each part of the New Testament is read in one step, from its texts to what they say. That step is the reading, and it is cited and rated, because texts do not entail conclusions and readings do. Each corpus is made to deliver only what its own texts say, and \emph{faith alone} is split into three parts, so that no corpus has to carry more than its share:

\begin{itemize}
  \item that faith is \textbf{necessary}: Paul (“not justified by works of the law except through faith”, Galatians 2:16; “the righteous shall live by faith”, 3:11), Hebrews (“without faith it is impossible to please God”) and John (“whoever does not believe is condemned already”);
  \item that faith is \textbf{sufficient}: Paul (“to the one who does not work but believes”, Romans 4:5), John, Peter and Jesus’ words in Luke;
  \item that works are \textbf{not the ground} of justification: Paul (Romans 3:28, Ephesians 2:9, Titus 3:5), Peter at the Jerusalem council, and Jesus’ parable of the Pharisee and the tax collector.
\end{itemize}
The rival is Trent, which holds every one of these texts. Set beside the canonical witness with James, it parts at exactly one reading: that “justified” in James 2:24 is the \emph{increase} of a justification already received, by works done in grace (Session VI, ch. 10).

\subsubsection*{What the early church said}
The readings are cited to the Reformers and to modern exegetes. They are also cited to fathers whom Rome itself venerates and cites, to show that the claims Trent condemns were not coined in the sixteenth century.

\begin{itemize}
  \item Clement of Rome: “not justified by ourselves … nor by works which we have wrought in holiness of heart; but by that faith through which, from the beginning, Almighty God has justified all men” (1 Clement 32.4).
  \item Polycarp: “by grace you are saved, not of works” (\emph{Philippians} 1).
  \item The Epistle to Diognetus, which the Catechism cites (CCC 2271): “that the righteousness of One should justify many transgressors” (9).
  \item Hilary of Poitiers, a Doctor of the Church: “faith alone justifies”, \emph{fides enim sola iustificat} (\emph{Commentary on Matthew} 8.6).
  \item Basil the Great, another: “to be justified by faith in Christ alone” (Homily 20, \emph{Of Humility}).
\end{itemize}
They show that faith’s sufficiency, and the exclusion of works as ground, were early Christian claims. They do not show that the fathers held the Reformers’ \emph{forensic} account of what justification is, a declaring righteous rather than a making righteous. Augustine reads justification as a making righteous, and Trent follows him there. Chrysostom’s seventh homily on Romans, on 3:27–28, speaks of justification by faith without saying “faith alone”. So the fathers are cited for what they say, and the account of what justification is is argued from Paul.

\subsubsection*{What this does not show}
It does not show that the Reformed readings are right. It shows where the disagreement is. Every text is common ground, every reading is cited to the people who hold it, and the results name, corpus by corpus, which readings a reader must reject to deny the conclusion.

\subsubsection*{The questions the later chapters take}
\begin{itemize}
  \item \textbf{What does each corpus say, and how is it read?} Each corpus’s texts and its one reading, with the ratings (“What each corpus says”).
  \item \textbf{What must an opponent reject, and does James contradict Paul?} No single corpus carries the case, and there are exactly three ways to overturn it. James coheres with Paul on one reading of 2:14, and Trent parts at one reading of 2:24 (“What follows, and what an opponent must reject”).
  \item \textbf{Where does each claim come from?} (“Appendix: the claims and their sources”).
\end{itemize}
\subsection*{What each corpus says}
\addcontentsline{toc}{subsection}{What each corpus says}
Each corpus is a line of reason: its texts as grounds, and one step — the reading — carrying them to what it delivers. The texts are granted by every party; the step is where each corpus is read, and it is cited and rated as a reading.

A corpus delivers only what its texts say. Hebrews says faith is necessary, and nothing about works; the Johannine texts say faith is sufficient and necessary; Peter and Jesus say that works are not the ground and that faith is sufficient. Paul says all three, and he is read twice, because he says them in different places: on works (Romans 3:28, Ephesians 2:8–9, Titus 3:5), that works are not the ground; on faith (Romans 4:4–5, Galatians 2:16, 3:11), that faith is necessary and sufficient. \emph{Faith alone} is the three together (\texttt{to\allowbreak{}Faith\allowbreak{}Alone}), so no reading has to carry more than its own texts.

\subsubsection*{The readings, rated}
Every reading but Hebrews’ is rated \texttt{disputed}: Trent grants the texts and denies that faith alone follows (Session VI, canon 9). Hebrews’ reading is \texttt{consensus}, because Trent makes it too — its chapter 8 calls faith the beginning of salvation, “without which it is impossible to please God”.

\subsubsection*{How this is checked}
\par\medskip\noindent\textbf{\texttt{pauline\allowbreak{}Reading\allowbreak{}Source}}\par\nobreak\noindent
Paul’s reading on works, cited to the Reformed exegesis, and to the first Christian writer outside the New Testament to say it: Clement of Rome, writing that we are justified “not by ourselves … nor by works which we have wrought in holiness of heart, but by that faith through which, from the beginning, Almighty God has justified all men” (1 Clement 32.4); and to Polycarp, who writes to the Philippians “knowing that by grace you are saved, not of works” (Philippians 1). Rome venerates both as saints.

{\small\begin{verbatim}
def paulineReadingSource : Source :=
  { primary := Reference.work schreinerFaithAlone,
    supporting :=
      [Reference.work Bib.calvinInstitutes (Bib.Locus.sectionRef "III.xi.19"),
        Reference.work anf1 (Bib.Locus.sectionRef "1 Clement 32.4"),
        Reference.work anf1
          (Bib.Locus.sectionRef "Polycarp, Epistle to the Philippians 1")],
    tradition := Tradition.reformedProtestant,
    confidence := Confidence.disputed }
\end{verbatim}}
\par\medskip\noindent\textbf{\texttt{pauline\allowbreak{}Faith\allowbreak{}Reading\allowbreak{}Source}}\par\nobreak\noindent
Paul’s reading on faith: that it is necessary — “not justified by works of the law except through faith” (Galatians 2:16, ἐὰν μή), “the righteous shall live by faith” (3:11) — and sufficient — “to the one who does not work but believes … his faith is counted as righteousness” (Romans 4:5). Trent makes the first half of it: its chapter 8 reads the Apostle’s “justified by faith” as faith “the beginning of human salvation, the foundation and root of all justification”. It denies the second (canon 9), so the reading is \texttt{disputed}.

Two fathers Rome names Doctors of the Church say the second half in terms. Basil: “recognizing oneself as lacking true justice, to be justified by faith in Christ alone” (Homily 20, \emph{Of Humility}). Hilary of Poitiers, of Jesus forgiving the paralytic: the Law could not absolve, “since faith alone justifies” — \emph{fides enim sola iustificat} (\emph{Commentary on Matthew} 8.6).

{\small\begin{verbatim}
def paulineFaithReadingSource : Source :=
  { primary := Reference.work schreinerFaithAlone,
    supporting :=
      [Reference.work tannerDecrees
          (Bib.Locus.sectionRef
            "Trent, Session VI (1547), Decree on Justification, ch. 8"),
        Reference.work Bib.calvinInstitutes (Bib.Locus.sectionRef "III.xi.19"),
        Reference.work basilAsceticalWorks
          (Bib.Locus.sectionRef "Homily 20, Of Humility"),
        Reference.work hilaryMatthew (Bib.Locus.sectionRef "8.6")],
    tradition := Tradition.reformedProtestant,
    confidence := Confidence.disputed }
\end{verbatim}}
\par\medskip\noindent\textbf{\texttt{hebrews\allowbreak{}Reading\allowbreak{}Source}}\par\nobreak\noindent
Hebrews’ reading: that faith is necessary. Conceded by Trent.

{\small\begin{verbatim}
def hebrewsReadingSource : Source :=
  {
    primary :=
      Reference.work bruceHebrews
        (Bib.Locus.adLoc { book := Book.hebrews, chapter := 11, verse := 6 }),
    supporting :=
      [Reference.work tannerDecrees
          (Bib.Locus.sectionRef
            "Trent, Session VI (1547), Decree on Justification, ch. 8")],
    tradition := Tradition.christianHistoricalGrammatical,
    confidence := Confidence.consensus }
\end{verbatim}}
\par\medskip\noindent\textbf{\texttt{johannine\allowbreak{}Reading\allowbreak{}Source}}\par\nobreak\noindent
The Johannine reading, cited to Calvin’s commentary on John. Aquinas reads the same texts of a faith formed by charity (\texttt{Testimony\allowbreak{}Arguments.Sola\allowbreak{}Fide}).

{\small\begin{verbatim}
def johannineReadingSource : Source :=
  {
    primary :=
      Reference.work calvinJohn
        (Bib.Locus.adLoc { book := Book.john, chapter := 3, verse := 16 }),
    tradition := Tradition.reformedProtestant,
    confidence := Confidence.disputed }
\end{verbatim}}
\par\medskip\noindent\textbf{\texttt{petrine\allowbreak{}Reading\allowbreak{}Source}}\par\nobreak\noindent
The Petrine reading, cited to Bruce on Peter’s speech at the Jerusalem council.

{\small\begin{verbatim}
def petrineReadingSource : Source :=
  {
    primary :=
      Reference.work bruceActs
        (Bib.Locus.adLoc { book := Book.acts, chapter := 15, verse := 11 }),
    tradition := Tradition.christianHistoricalGrammatical,
    confidence := Confidence.disputed }
\end{verbatim}}
\par\medskip\noindent\textbf{\texttt{dominical\allowbreak{}Reading\allowbreak{}Source}}\par\nobreak\noindent
The dominical reading, cited to Marshall on the Pharisee and the tax collector.

{\small\begin{verbatim}
def dominicalReadingSource : Source :=
  {
    primary :=
      Reference.work marshallLuke
        (Bib.Locus.adLoc { book := Book.luke, chapter := 18, verse := 14 }),
    tradition := Tradition.christianHistoricalGrammatical,
    confidence := Confidence.disputed }
\end{verbatim}}
\par\medskip\noindent\textbf{\texttt{james\allowbreak{}Reading\allowbreak{}Source}}\par\nobreak\noindent
James’s reading: works as the fruit of a faith that has justified.

{\small\begin{verbatim}
def jamesReadingSource : Source :=
  { primary := Reference.work mooJames (Bib.Locus.adLoc Scripture.james2_14),
    supporting :=
      [Reference.work johnsonJames (Bib.Locus.adLoc Scripture.james2_24),
        Reference.work Bib.calvinInstitutes
          (Bib.Locus.sectionRef "III.xvii.11")],
    tradition := Tradition.reformedProtestant,
    confidence := Confidence.disputed }
\end{verbatim}}
\par\medskip\noindent\textbf{\texttt{trent\allowbreak{}Against\allowbreak{}Sufficiency\allowbreak{}Source}}\par\nobreak\noindent
Trent’s use of James 2:24 against faith alone: the verse it quotes in chapter 10, against the sufficiency canon 9 denies.

{\small\begin{verbatim}
def trentAgainstSufficiencySource : Source :=
  {
    primary :=
      Reference.work tannerDecrees
        (Bib.Locus.sectionRef
          "Trent, Session VI (1547), Decree on Justification, ch. 10 and canon 9"),
    tradition := Tradition.romanCatholic, confidence := Confidence.disputed }
\end{verbatim}}
\par\medskip\noindent\textbf{\texttt{trent\allowbreak{}On\allowbreak{}James\allowbreak{}Source}}\par\nobreak\noindent
Trent’s reading of James 2:24: justification increased by works.

{\small\begin{verbatim}
def trentOnJamesSource : Source :=
  {
    primary :=
      Reference.work tannerDecrees
        (Bib.Locus.sectionRef
          "Trent, Session VI (1547), Decree on Justification, ch. 10"),
    tradition := Tradition.romanCatholic, confidence := Confidence.disputed }
\end{verbatim}}
\paragraph*{The lines}
\leavevmode\par\nobreak
\par\medskip\noindent\textbf{\texttt{pauline\allowbreak{}Line}} --- Paul, on works (Romans 3:28, Ephesians 2:8–9, Titus 3:5)\par\nobreak\noindent
\textbf{Paul, on works}, read as canon: the undisputed letters with Ephesians and Titus, whoever wrote them. “Apart from works of the law” (Romans 3:28), “not of works” (Ephesians 2:9), “not by works done in righteousness” (Titus 3:5): works are not the ground.

\begin{derivation}
  \item[(1)] \(\pvar{TestimonyArguments.CanonicalWitness}{1}\)\gloss{Romans 3:28: a person is justified by faith apart from works of the law}
  \item[(2)] \(\pvar{TestimonyArguments.CanonicalWitness}{5}\)\gloss{Ephesians 2:8–9: saved by grace through faith, not of works}
  \item[(3)] \(\pvar{TestimonyArguments.CanonicalWitness}{6}\)\gloss{Titus 3:5: saved not by works done in righteousness but by his mercy}
  \item[(4)] \((\pvar{TestimonyArguments.CanonicalWitness}{1} \land \pvar{TestimonyArguments.CanonicalWitness}{5} \land \pvar{TestimonyArguments.CanonicalWitness}{6}) \rightarrow \pvar{TestimonyArguments.CanonicalWitness}{23}\)
  \item[\(\vdash\)] \(\pvar{TestimonyArguments.CanonicalWitness}{23}\)\gloss{\emph{at issue:} Works are not the ground of justification}
\end{derivation}
\par\medskip\noindent\textbf{\texttt{pauline\allowbreak{}Faith\allowbreak{}Line}} --- Paul, on faith (Romans 4:4–5, Galatians 2:16, 3:11)\par\nobreak\noindent
\textbf{Paul, on faith}: faith is necessary — no one is justified except through it (Galatians 2:16, 3:11) — and sufficient — to the one who does not work but believes, faith is counted as righteousness (Romans 4:4–5). Paul’s own message, stated apart from what he says about works.

\begin{derivation}
  \item[(1)] \(\pvar{TestimonyArguments.CanonicalWitness}{2}\)\gloss{Romans 4:4–5: to the one who does not work but believes, faith is counted as righteousness}
  \item[(2)] \(\pvar{TestimonyArguments.CanonicalWitness}{3}\)\gloss{Galatians 2:16: not justified by works of the law but through faith in Christ}
  \item[(3)] \(\pvar{TestimonyArguments.CanonicalWitness}{4}\)\gloss{Galatians 3:11: no one is justified before God by the law; the righteous shall live by faith}
  \item[(4)] \((\pvar{TestimonyArguments.CanonicalWitness}{2} \land \pvar{TestimonyArguments.CanonicalWitness}{3} \land \pvar{TestimonyArguments.CanonicalWitness}{4}) \rightarrow \pvar{TestimonyArguments.CanonicalWitness}{21} \land \pvar{TestimonyArguments.CanonicalWitness}{22}\)
  \item[\(\vdash\)] \(\pvar{TestimonyArguments.CanonicalWitness}{21} \land \pvar{TestimonyArguments.CanonicalWitness}{22}\)\gloss{\emph{at issue:} \(\pvar{TestimonyArguments.CanonicalWitness}{21}\) Faith is necessary: no one is saved without it; \(\pvar{TestimonyArguments.CanonicalWitness}{22}\) Faith is sufficient: whoever believes is saved}
\end{derivation}
\par\medskip\noindent\textbf{\texttt{hebrews\allowbreak{}Line}} --- Hebrews\par\nobreak\noindent
\textbf{Hebrews}, read as its own corpus: faith is necessary.

\begin{derivation}
  \item[(1)] \(\pvar{TestimonyArguments.CanonicalWitness}{7}\)\gloss{Hebrews 10:38–39: the righteous one shall live by faith}
  \item[(2)] \(\pvar{TestimonyArguments.CanonicalWitness}{8}\)\gloss{Hebrews 11:6: without faith it is impossible to please God}
  \item[(3)] \((\pvar{TestimonyArguments.CanonicalWitness}{7} \land \pvar{TestimonyArguments.CanonicalWitness}{8}) \rightarrow \pvar{TestimonyArguments.CanonicalWitness}{21}\)
  \item[\(\vdash\)] \(\pvar{TestimonyArguments.CanonicalWitness}{21}\)\gloss{\emph{at issue:} Faith is necessary: no one is saved without it}
\end{derivation}
\par\medskip\noindent\textbf{\texttt{johannine\allowbreak{}Line}} --- John\par\nobreak\noindent
\textbf{John}: whoever believes has life, and whoever does not is condemned already.

\begin{derivation}
  \item[(1)] \(\pvar{TestimonyArguments.CanonicalWitness}{9}\)\gloss{John 3:16–18: whoever believes has eternal life; whoever does not believe is condemned already}
  \item[(2)] \(\pvar{TestimonyArguments.CanonicalWitness}{10}\)\gloss{John 5:24: whoever hears and believes has eternal life and does not come into judgement}
  \item[(3)] \(\pvar{TestimonyArguments.CanonicalWitness}{11}\)\gloss{John 6:28–29: the work of God is that you believe in him whom he has sent}
  \item[(4)] \((\pvar{TestimonyArguments.CanonicalWitness}{9} \land \pvar{TestimonyArguments.CanonicalWitness}{10} \land \pvar{TestimonyArguments.CanonicalWitness}{11}) \rightarrow \pvar{TestimonyArguments.CanonicalWitness}{22} \land \pvar{TestimonyArguments.CanonicalWitness}{21}\)
  \item[\(\vdash\)] \(\pvar{TestimonyArguments.CanonicalWitness}{22} \land \pvar{TestimonyArguments.CanonicalWitness}{21}\)\gloss{\emph{at issue:} \(\pvar{TestimonyArguments.CanonicalWitness}{22}\) Faith is sufficient: whoever believes is saved; \(\pvar{TestimonyArguments.CanonicalWitness}{21}\) Faith is necessary: no one is saved without it}
\end{derivation}
\par\medskip\noindent\textbf{\texttt{petrine\allowbreak{}Line}} --- Peter (Acts 10 and 15)\par\nobreak\noindent
\textbf{Peter}, in Luke’s report: forgiveness to everyone who believes, and salvation through grace, without the yoke.

\begin{derivation}
  \item[(1)] \(\pvar{TestimonyArguments.CanonicalWitness}{12}\)\gloss{Acts 10:43: everyone who believes in him receives forgiveness of sins}
  \item[(2)] \(\pvar{TestimonyArguments.CanonicalWitness}{13}\)\gloss{Acts 15:10–11: no yoke on the disciples; we are saved through the grace of the Lord Jesus}
  \item[(3)] \((\pvar{TestimonyArguments.CanonicalWitness}{12} \land \pvar{TestimonyArguments.CanonicalWitness}{13}) \rightarrow \pvar{TestimonyArguments.CanonicalWitness}{22} \land \pvar{TestimonyArguments.CanonicalWitness}{23}\)
  \item[\(\vdash\)] \(\pvar{TestimonyArguments.CanonicalWitness}{22} \land \pvar{TestimonyArguments.CanonicalWitness}{23}\)\gloss{\emph{at issue:} \(\pvar{TestimonyArguments.CanonicalWitness}{22}\) Faith is sufficient: whoever believes is saved; \(\pvar{TestimonyArguments.CanonicalWitness}{23}\) Works are not the ground of justification}
\end{derivation}
\par\medskip\noindent\textbf{\texttt{dominical\allowbreak{}Line}} --- Jesus (Luke 7 and 18)\par\nobreak\noindent
\textbf{Jesus}, in Luke: the woman whose faith saved her, and the tax collector who went home justified.

\begin{derivation}
  \item[(1)] \(\pvar{TestimonyArguments.CanonicalWitness}{14}\)\gloss{Luke 7:50: your faith has saved you}
  \item[(2)] \(\pvar{TestimonyArguments.CanonicalWitness}{15}\)\gloss{Luke 18:9–14: the tax collector, not the Pharisee who listed his works, went home justified}
  \item[(3)] \((\pvar{TestimonyArguments.CanonicalWitness}{14} \land \pvar{TestimonyArguments.CanonicalWitness}{15}) \rightarrow \pvar{TestimonyArguments.CanonicalWitness}{22} \land \pvar{TestimonyArguments.CanonicalWitness}{23}\)
  \item[\(\vdash\)] \(\pvar{TestimonyArguments.CanonicalWitness}{22} \land \pvar{TestimonyArguments.CanonicalWitness}{23}\)\gloss{\emph{at issue:} \(\pvar{TestimonyArguments.CanonicalWitness}{22}\) Faith is sufficient: whoever believes is saved; \(\pvar{TestimonyArguments.CanonicalWitness}{23}\) Works are not the ground of justification}
\end{derivation}
\par\medskip\noindent\textbf{\texttt{to\allowbreak{}Faith\allowbreak{}Alone}}\par\nobreak\noindent
The three parts, together, are \emph{faith alone}.

\begin{derivation}
  \item[] \((\pvar{TestimonyArguments.CanonicalWitness}{21} \land \pvar{TestimonyArguments.CanonicalWitness}{22} \land \pvar{TestimonyArguments.CanonicalWitness}{23}) \rightarrow \pvar{TestimonyArguments.CanonicalWitness}{24}\)\gloss{\emph{would give:} Justification is by faith alone}
\end{derivation}
\par\medskip\noindent\textbf{\texttt{james\allowbreak{}Line}} --- James 2\par\nobreak\noindent
\textbf{James}, read with the distinction between justification and sanctification: the faith he calls dead is the demons’ assent, and the justification he means is the display of a faith already counted as righteousness. So works are the fruit of saving faith, not its ground.

\begin{derivation}
  \item[(1)] \(\pvar{TestimonyArguments.CanonicalWitness}{16}\)\gloss{James 2:14–17: faith by itself, if it does not have works, is dead}
  \item[(2)] \(\pvar{TestimonyArguments.CanonicalWitness}{17}\)\gloss{James 2:19: even the demons believe, and shudder}
  \item[(3)] \(\pvar{TestimonyArguments.CanonicalWitness}{18}\)\gloss{James 2:21–23: Abraham’s faith was completed by his works, and Genesis 15:6 was fulfilled}
  \breakrow{(4)} \(\pvar{TestimonyArguments.CanonicalWitness}{19}\)\gloss{James 2:24: a person is justified by works and not by faith alone}
  \breakrow{(5)} \(\pvar{TestimonyArguments.CanonicalWitness}{27}\)\gloss{James 2:14–26 targets a barren faith — mere assent — not Paul’s doctrine}
  \breakrow{(6)} \(\pvar{TestimonyArguments.CanonicalWitness}{28}\)\gloss{‘Justified’ in James 2:21–24 means shown to be righteous (denied by Trent, which reads 2:24 as an increase of justification by works done in grace)}
  \breakrow{(7)} \(\pvar{TestimonyArguments.CanonicalWitness}{29}\)\gloss{Justification and sanctification are inseparable but distinct}
  \item[(8)] \((\pvar{TestimonyArguments.CanonicalWitness}{16} \land \pvar{TestimonyArguments.CanonicalWitness}{18} \land \pvar{TestimonyArguments.CanonicalWitness}{27} \land \pvar{TestimonyArguments.CanonicalWitness}{28} \land \pvar{TestimonyArguments.CanonicalWitness}{29}) \rightarrow \pvar{TestimonyArguments.CanonicalWitness}{25}\)
  \item[\(\vdash\)] \(\pvar{TestimonyArguments.CanonicalWitness}{25}\)\gloss{\emph{at issue:} Good works are the fruit and evidence of a true and lively faith}
\end{derivation}
\par\medskip\noindent\textbf{\texttt{to\allowbreak{}Never\allowbreak{}Alone}}\par\nobreak\noindent
Faith alone, with its fruit, is faith that is never alone.

\begin{derivation}
  \item[] \((\pvar{TestimonyArguments.CanonicalWitness}{24} \land \pvar{TestimonyArguments.CanonicalWitness}{25}) \rightarrow \pvar{TestimonyArguments.CanonicalWitness}{26}\)\gloss{\emph{would give:} Justification is by faith alone, by a faith that is never alone}
\end{derivation}
\par\medskip\noindent\textbf{\texttt{trent\allowbreak{}On\allowbreak{}James}}\par\nobreak\noindent
\textbf{Trent’s reading of James 2:24.} “Justified” there is the increase of a justification already received, by works done in grace; so works are part of its ground, and faith alone is denied.

\begin{derivation}
  \item[] \((\pvar{TestimonyArguments.CanonicalWitness}{19} \land \lnot \pvar{TestimonyArguments.CanonicalWitness}{28}) \rightarrow \lnot \pvar{TestimonyArguments.CanonicalWitness}{23}\)\gloss{\emph{would deny:} Works are not the ground of justification}
\end{derivation}
\par\medskip\noindent\textbf{\texttt{james\allowbreak{}Against\allowbreak{}Sufficiency}}\par\nobreak\noindent
\textbf{Trent’s reading of James against sufficiency.} James says faith without works is dead and cannot save (2:14–17), and that a person is justified “not by faith alone” (2:24); read of faith as such, that denies that faith suffices. The rival to the reading on which James and Paul cohere: there, the faith that cannot save is the faith 2:14–17 describes, faith without works (\texttt{james2\allowbreak{}Targets\allowbreak{}Dead\allowbreak{}Faith}).

\begin{derivation}
  \item[] \((\pvar{TestimonyArguments.CanonicalWitness}{16} \land \pvar{TestimonyArguments.CanonicalWitness}{19}) \rightarrow \lnot \pvar{TestimonyArguments.CanonicalWitness}{22}\)\gloss{\emph{would deny:} Faith is sufficient: whoever believes is saved}
\end{derivation}
\par\medskip\noindent\textbf{\texttt{not\allowbreak{}Alone\allowbreak{}If\allowbreak{}Works\allowbreak{}Ground}}\par\nobreak\noindent
And what follows from that: not faith alone.

\begin{derivation}
  \item[] \(\lnot \pvar{TestimonyArguments.CanonicalWitness}{23} \rightarrow \lnot \pvar{TestimonyArguments.CanonicalWitness}{24}\)\gloss{\emph{would deny:} Justification is by faith alone}
\end{derivation}
\par\medskip\noindent\textbf{\texttt{corpora}}\par\nobreak\noindent
The corpora, one reading each, Paul’s two readings first.

{\small\begin{verbatim}
def corpora : List (Line Claim) :=
  [paulineLine, paulineFaithLine,
    hebrewsLine, johannineLine,
    petrineLine, dominicalLine]
\end{verbatim}}
\paragraph*{The readings, by name}
\leavevmode\par\nobreak
What an opponent can reject is a corpus’s reading. The burden numbers the readings by their place among the lines — the corpora, then James — and these names say which is which, so every burden below reads as the corpora it concerns. \texttt{readings\_\allowbreak{}named} checks each name against its line.

\par\medskip\noindent\textbf{\texttt{readings\_\allowbreak{}named}}\par\nobreak\noindent
Each reading’s name points at its own line.

{\small\begin{verbatim}
theorem readings_named : (corpora ++ [jamesLine])[Reading.paul]? = some
    paulineLine ∧ (corpora ++ [jamesLine])[Reading.paulOnFaith]? = some
    paulineFaithLine ∧ (corpora ++ [jamesLine])[Reading.hebrews]? = some
    hebrewsLine ∧ (corpora ++ [jamesLine])[Reading.john]? = some johannineLine
    ∧ (corpora ++ [jamesLine])[Reading.peter]? = some petrineLine ∧ (corpora
    ++ [jamesLine])[Reading.jesus]? = some dominicalLine ∧ (corpora ++
    [jamesLine])[Reading.james]? = some jamesLine
-- axioms: propext
\end{verbatim}}
\subsection*{What follows, and what an opponent must reject}
\addcontentsline{toc}{subsection}{What follows, and what an opponent must reject}
Grant the texts, read each corpus as its line reads it, and faith alone follows (\texttt{canonical\_\allowbreak{}establishes}). Trent’s position, which holds the same texts and reads James 2:24 of an increase of justification, follows from its own premises too (\texttt{tridentine\_\allowbreak{}establishes}). Every result here is conditional on its premises, as everywhere in the library. What this argument adds is that every premise but the readings is a text all parties accept as Scripture.

\subsubsection*{What an opponent must reject}
Because the texts are granted, what an opponent can reject is a reading. The \emph{opponent’s burden} lists every smallest set of readings whose rejection overturns the case, and proves there is no other way (\texttt{canonical\_\allowbreak{}burden}). The library computes the sets, and no one picks them. There are exactly three, one for each part of faith alone.

\begin{itemize}
  \item To deny that works are not the ground, an opponent must reject three readings: Paul’s on works, Peter’s and Jesus’. Trent’s position, as encoded here, takes this way: on its reading of James 2:24, works are part of the ground.
  \item To deny that faith is necessary, three: Paul’s on faith, Hebrews’ and John’s. Trent does not take this way, since its chapter 8 reads the Apostle’s “justified by faith” of faith as “the beginning of human salvation … without which it is impossible to please God”.
  \item To deny that faith suffices, four: Paul’s on faith, John’s, Peter’s and Jesus’. Trent’s canon 9, which condemns the claim that the ungodly is justified by faith alone, denies this too. The reading of James that goes with it is set out below as a rival of its own, not as part of Trent’s position.
\end{itemize}
Hebrews’ reading is rated \texttt{consensus}, the highest rating, because Trent makes it too, and every other reading \texttt{disputed}, the lowest. Because every set holds more than one corpus, no single corpus is load-bearing: take any one reading away and the conclusion still follows (\texttt{paul\_\allowbreak{}not\_\allowbreak{}load\_\allowbreak{}bearing}, \texttt{hebrews\_\allowbreak{}not\_\allowbreak{}load\_\allowbreak{}bearing}, \texttt{john\_\allowbreak{}not\_\allowbreak{}load\_\allowbreak{}bearing}, \texttt{peter\_\allowbreak{}not\_\allowbreak{}load\_\allowbreak{}bearing}, \texttt{jesus\_\allowbreak{}not\_\allowbreak{}load\_\allowbreak{}bearing}). Two corpora are enough, Hebrews and Jesus’ words (\texttt{hebrews\_\allowbreak{}and\_\allowbreak{}jesus\_\allowbreak{}suffice}).

The burden is how this library lets several witnesses matter without counting them. When three corpora say works are not the ground, an opponent has to break three readings, each on its own texts.

\subsubsection*{Paul and James}
Paul’s message about faith can be stated on its own: faith is necessary and it is sufficient, from Romans 4:4–5 and Galatians 2:16 and 3:11 (\texttt{paul\_\allowbreak{}on\_\allowbreak{}faith\_\allowbreak{}establishes}). Does James contradict it? Not on its texts, since none of them says either. It turns on one reading, of whom James means when he asks “can that faith save him?” (2:14). Read 2:14–17 as about faith without works, and leave aside the contested sense of “justified”, and Paul can be held together with every text of James (\texttt{paul\_\allowbreak{}on\_\allowbreak{}faith\_\allowbreak{}stands\_\allowbreak{}with\_\allowbreak{}james}). James does not contradict Paul on that reading, and does not support him either. Read James’s “that faith” (2:14) and “not by faith alone” (2:24) of faith as such, as Trent does, and the two cannot both be held (\texttt{james\_\allowbreak{}against\_\allowbreak{}sufficiency\_\allowbreak{}contradicts\_\allowbreak{}paul}). So the harmony depends on how 2:14 is read.

James does not argue for faith alone, and he does not argue against works. He argues against a faith without them: the demons’ belief of 2:19, and the faith that says “go in peace” and gives nothing (2:14–17). Reformed readers distinguish justification, God’s verdict, from sanctification, the believer’s renewal. Read with that distinction, James 2:21–24 uses “justified” in a \emph{demonstrative} sense: Abraham’s offering of Isaac \emph{showed} the faith already counted as righteousness in Genesis 15:6, which James himself quotes. So James adds what the other corpora do not say, that works are the fruit of saving faith. With him the conclusion becomes the Reformed formula, justification by faith alone, by a faith that is never alone. That is a paraphrase of the Westminster Confession, which calls faith “the alone instrument of justification: yet is it not alone in the person justified” (XI.2). Every text of James is a premise, 2:24 included. Jude, which does not address the question, is granted as consistent with the rest and nothing more.

\subsubsection*{Where Trent parts}
Set beside the canonical witness with James, Trent parts at that sense of “justified”. On Trent’s reading of James 2:24, works are part of the ground, and faith alone fails. Trent’s own decree says more than its reading of James. Its chapter 8 holds that nothing before justification, “whether faith or works”, merits it, so Trent too denies that works ground the first justification. Its canon 24 condemns calling works “merely the fruits and signs of Justification obtained”, which answers the demonstrative reading directly. The encoding takes Trent’s position at chapter 10 only. The crux is found mechanically, as the one claim that breaks Trent at its reading of James (\texttt{why\allowbreak{}The\allowbreak{}Canonical\allowbreak{}Witness\allowbreak{}Stands\allowbreak{}Against\allowbreak{}Trent}). The Reformed formula needs it: without the demonstrative reading of James, “faith never alone” no longer follows (\texttt{james\_\allowbreak{}demonstrative\_\allowbreak{}is\_\allowbreak{}load\_\allowbreak{}bearing}). The burden with James says the same at the level of readings. Rejecting James’s reading alone overturns the formula, and nothing else about faith alone changes (\texttt{canonical\allowbreak{}With\allowbreak{}James\_\allowbreak{}burden}).

\subsubsection*{The positions}
Trent’s position is written first, because the rival is written before anything is proved. Then come the canonical case, the case with James, and the particular burdens a reader can check one at a time: the case with one corpus’s reading rejected, with the readings behind one part of faith alone rejected together, and with all but two rejected. Every text stays granted in each, and what goes is the step that reads them.

\subsubsection*{How this is checked}
\paragraph*{The rival}
\leavevmode\par\nobreak
\par\medskip\noindent\textbf{\texttt{tridentine\allowbreak{}Reading}}\par\nobreak\noindent
Trent’s world: every text holds, “justified” in James 2:24 is not demonstrative, works are part of the ground, and faith alone fails.

{\small\begin{verbatim}
def tridentineReading : Valuation Claim :=
  fun a =>
    match a with
    | Claim.jamesJustifiesDemonstratively => False
    | Claim.worksAreNotTheGround => False
    | Claim.justificationByFaithAlone => False
    | x => True
\end{verbatim}}
\par\medskip\noindent\textbf{\texttt{tridentine\_\allowbreak{}establishes}}\par\nobreak\noindent
\textbf{Trent’s reading is coherent, and it denies faith alone.} Granting every text, Trent reaches its conclusion by its reading of James 2:24.

{\small\begin{verbatim}
theorem tridentine_establishes : Establishes tridentine
-- axioms: propext, Quot.sound
\end{verbatim}}
\par\medskip\noindent\textbf{\texttt{tridentine\_\allowbreak{}is\_\allowbreak{}satisfiable}}\par\nobreak\noindent
Trent’s premises have a model.

{\small\begin{verbatim}
theorem tridentine_is_satisfiable : Satisfiable tridentine.premises
-- axioms: propext, Quot.sound
\end{verbatim}}
\paragraph*{The canonical witness}
\leavevmode\par\nobreak
\par\medskip\noindent\textbf{\texttt{canonical\allowbreak{}Reading}}\par\nobreak\noindent
The world in which every text and every reading holds.

{\small\begin{verbatim}
def canonicalReading : Valuation Claim :=
  fun x => True
\end{verbatim}}
\par\medskip\noindent\textbf{\texttt{canonical\_\allowbreak{}establishes}}\par\nobreak\noindent
\textbf{From the texts all parties accept, faith alone follows}, on one reading per corpus: Paul, Hebrews, John, Peter and Jesus’ words, each delivering only what its own texts say.

{\small\begin{verbatim}
theorem canonical_establishes : Establishes canonicalCase
-- axioms: propext, Quot.sound
\end{verbatim}}
\par\medskip\noindent\textbf{\texttt{canonical\allowbreak{}Case\_\allowbreak{}is\_\allowbreak{}satisfiable}}\par\nobreak\noindent
The canonical case has a model.

{\small\begin{verbatim}
theorem canonicalCase_is_satisfiable : Satisfiable canonicalCase.premises
-- axioms: propext, Quot.sound
\end{verbatim}}
\par\medskip\noindent\textbf{\texttt{canonical\allowbreak{}With\allowbreak{}James\_\allowbreak{}establishes}}\par\nobreak\noindent
\textbf{With James, the Reformed formula follows}: justification by faith alone, by a faith that is never alone. James’s texts — all of them, 2:24 included — are premises, and they cohere with the other corpora on the distinction between justification and sanctification.

{\small\begin{verbatim}
theorem canonicalWithJames_establishes : Establishes canonicalWithJames
-- axioms: propext, Quot.sound
\end{verbatim}}
\par\medskip\noindent\textbf{\texttt{canonical\allowbreak{}With\allowbreak{}James\_\allowbreak{}is\_\allowbreak{}satisfiable}}\par\nobreak\noindent
The case with James has a model: James and the other corpora can all be held at once.

{\small\begin{verbatim}
theorem canonicalWithJames_is_satisfiable : Satisfiable
    canonicalWithJames.premises
-- axioms: propext, Quot.sound
\end{verbatim}}
\paragraph*{What an opponent must reject}
\leavevmode\par\nobreak
The texts are granted, so what an opponent can reject is a reading. The burden lists every minimal set of readings whose rejection overturns the case, and proves there is no other way (\texttt{Testimony.Logic.Burden}). The results after it take its sets one at a time, as named packages a reader can check: each is an instance of the burden, and is proved from it.

\par\medskip\noindent\textbf{\texttt{canonical\_\allowbreak{}burden}}\par\nobreak\noindent
\textbf{The opponent’s burden against the canonical witness.} There are exactly three ways to overturn it, one for each part of faith alone:

\begin{itemize}
  \item reject the three readings that say works are not the ground — Paul’s on works, Peter’s and Jesus’;
  \item or the three that say faith is necessary — Paul’s on faith, Hebrews’ and John’s;
  \item or the four that say faith is sufficient — Paul’s on faith, John’s, Peter’s and Jesus’.
\end{itemize}
Nothing less overturns it, and nothing else does. Trent does not take the second way: its Decree on Justification reads the Apostle’s “justified by faith” as faith “the beginning of human salvation … without which it is impossible to please God”.

{\small\begin{verbatim}
theorem canonical_burden : OpponentsBurden canonicalCase corpora []
    [toFaithAlone] [[Reading.paul, Reading.peter, Reading.jesus],
    [Reading.paulOnFaith, Reading.hebrews, Reading.john],
    [Reading.paulOnFaith, Reading.john, Reading.peter, Reading.jesus]]
-- axioms: propext, Classical.choice, Quot.sound
\end{verbatim}}
\paragraph{What an opponent must reject.} \emph{The canonical witness to faith alone} concludes that justification is by faith alone. That conclusion is overturned exactly when every reading in one of these sets is rejected. Each set is minimal, and no other rejection overturns it:
\begin{enumerate}
\item \emph{Paul, on works (Romans 3:28, Ephesians 2:8–9, Titus 3:5)} (disputed); \emph{Peter (Acts 10 and 15)} (disputed); \emph{Jesus (Luke 7 and 18)} (disputed)
\item \emph{Paul, on faith (Romans 4:4–5, Galatians 2:16, 3:11)} (disputed); \emph{Hebrews} (consensus); \emph{John} (disputed)
\item \emph{Paul, on faith (Romans 4:4–5, Galatians 2:16, 3:11)} (disputed); \emph{John} (disputed); \emph{Peter (Acts 10 and 15)} (disputed); \emph{Jesus (Luke 7 and 18)} (disputed)
\end{enumerate}
\paragraph*{No single corpus carries it}
\leavevmode\par\nobreak
Every set in the burden holds more than one corpus, so no single corpus’s rejection overturns the case. Here is each, taken away in turn.

\par\medskip\noindent\textbf{\texttt{paul\_\allowbreak{}not\_\allowbreak{}load\_\allowbreak{}bearing}}\par\nobreak\noindent
\textbf{Paul is not load-bearing.} Reject both his readings, and Peter and Jesus still say that works are not the ground, Hebrews and John that faith is necessary, and John, Peter and Jesus that it suffices.

{\small\begin{verbatim}
theorem paul_not_load_bearing : Establishes withoutPaul
-- axioms: propext, Classical.choice, Quot.sound
\end{verbatim}}
\par\medskip\noindent\textbf{\texttt{hebrews\_\allowbreak{}not\_\allowbreak{}load\_\allowbreak{}bearing}}\par\nobreak\noindent
\textbf{Nor is Hebrews.} Paul and John say faith is necessary too.

{\small\begin{verbatim}
theorem hebrews_not_load_bearing : Establishes withoutHebrews
-- axioms: propext, Classical.choice, Quot.sound
\end{verbatim}}
\par\medskip\noindent\textbf{\texttt{john\_\allowbreak{}not\_\allowbreak{}load\_\allowbreak{}bearing}}\par\nobreak\noindent
\textbf{Nor is John.} Paul and Hebrews say faith is necessary, and Paul, Peter and Jesus that it is sufficient.

{\small\begin{verbatim}
theorem john_not_load_bearing : Establishes withoutJohn
-- axioms: propext, Classical.choice, Quot.sound
\end{verbatim}}
\par\medskip\noindent\textbf{\texttt{peter\_\allowbreak{}not\_\allowbreak{}load\_\allowbreak{}bearing}}\par\nobreak\noindent
\textbf{Nor is Peter.} Paul and Jesus still say that works are not the ground.

{\small\begin{verbatim}
theorem peter_not_load_bearing : Establishes withoutPeter
-- axioms: propext, Classical.choice, Quot.sound
\end{verbatim}}
\par\medskip\noindent\textbf{\texttt{jesus\_\allowbreak{}not\_\allowbreak{}load\_\allowbreak{}bearing}}\par\nobreak\noindent
\textbf{Nor are Jesus’ words.} Paul and Peter still say that works are not the ground.

{\small\begin{verbatim}
theorem jesus_not_load_bearing : Establishes withoutJesus
-- axioms: propext, Classical.choice, Quot.sound
\end{verbatim}}
\par\medskip\noindent\textbf{\texttt{hebrews\_\allowbreak{}and\_\allowbreak{}jesus\_\allowbreak{}suffice}}\par\nobreak\noindent
\textbf{Two corpora are enough}: Hebrews, that faith is necessary, and Jesus’ words, that it is sufficient and that works are not the ground. Rejecting Paul’s, John’s and Peter’s readings leaves a reading of each set standing.

{\small\begin{verbatim}
theorem hebrews_and_jesus_suffice : Establishes hebrewsAndJesus
-- axioms: propext, Classical.choice, Quot.sound
\end{verbatim}}
\par\medskip\noindent\textbf{\texttt{without\allowbreak{}Paul\_\allowbreak{}is\_\allowbreak{}satisfiable}}\par\nobreak\noindent
Each named variant has a model: rejecting readings never makes the rest inconsistent. Without Paul:

{\small\begin{verbatim}
theorem withoutPaul_is_satisfiable : Satisfiable withoutPaul.premises
-- axioms: propext, Classical.choice, Quot.sound
\end{verbatim}}
\par\medskip\noindent\textbf{\texttt{without\allowbreak{}Hebrews\_\allowbreak{}is\_\allowbreak{}satisfiable}}\par\nobreak\noindent
Without Hebrews, the premises have a model.

{\small\begin{verbatim}
theorem withoutHebrews_is_satisfiable : Satisfiable withoutHebrews.premises
-- axioms: propext, Classical.choice, Quot.sound
\end{verbatim}}
\par\medskip\noindent\textbf{\texttt{without\allowbreak{}John\_\allowbreak{}is\_\allowbreak{}satisfiable}}\par\nobreak\noindent
Without John, the premises have a model.

{\small\begin{verbatim}
theorem withoutJohn_is_satisfiable : Satisfiable withoutJohn.premises
-- axioms: propext, Classical.choice, Quot.sound
\end{verbatim}}
\par\medskip\noindent\textbf{\texttt{without\allowbreak{}Peter\_\allowbreak{}is\_\allowbreak{}satisfiable}}\par\nobreak\noindent
Without Peter, the premises have a model.

{\small\begin{verbatim}
theorem withoutPeter_is_satisfiable : Satisfiable withoutPeter.premises
-- axioms: propext, Classical.choice, Quot.sound
\end{verbatim}}
\par\medskip\noindent\textbf{\texttt{without\allowbreak{}Jesus\_\allowbreak{}is\_\allowbreak{}satisfiable}}\par\nobreak\noindent
Without Jesus’ words, the premises have a model.

{\small\begin{verbatim}
theorem withoutJesus_is_satisfiable : Satisfiable withoutJesus.premises
-- axioms: propext, Classical.choice, Quot.sound
\end{verbatim}}
\par\medskip\noindent\textbf{\texttt{hebrews\allowbreak{}And\allowbreak{}Jesus\_\allowbreak{}is\_\allowbreak{}satisfiable}}\par\nobreak\noindent
Hebrews and Jesus’ words alone have a model.

{\small\begin{verbatim}
theorem hebrewsAndJesus_is_satisfiable : Satisfiable hebrewsAndJesus.premises
-- axioms: propext, Classical.choice, Quot.sound
\end{verbatim}}
\paragraph*{The three ways, each written down}
\leavevmode\par\nobreak
Each way of overturning the case is a world an opponent could stand in: every text true, the readings of the other corpora granted, and faith alone false. Naming that world is what makes the burden concrete.

\par\medskip\noindent\textbf{\texttt{no\allowbreak{}Works\allowbreak{}Witness\allowbreak{}Reading}}\par\nobreak\noindent
The first way: grant every text and every reading on faith, but none of Paul’s on works, Peter’s or Jesus’. Faith is necessary and sufficient, but works are not shown not to be the ground.

{\small\begin{verbatim}
def noWorksWitnessReading : Valuation Claim :=
  fun a =>
    match a with
    | Claim.worksAreNotTheGround => False
    | Claim.justificationByFaithAlone => False
    | x => True
\end{verbatim}}
\par\medskip\noindent\textbf{\texttt{works\_\allowbreak{}witnesses\_\allowbreak{}jointly\_\allowbreak{}load\_\allowbreak{}bearing}}\par\nobreak\noindent
\textbf{To deny that works are not the ground, an opponent must reject three readings} — Paul’s on works, Peter’s and Jesus’. With all three rejected, faith alone no longer follows; with any one kept, it does.

{\small\begin{verbatim}
theorem works_witnesses_jointly_load_bearing : ¬Establishes
    withoutTheWorksWitnesses
-- axioms: propext, Quot.sound
\end{verbatim}}
\par\medskip\noindent\textbf{\texttt{no\allowbreak{}Necessity\allowbreak{}Witness\allowbreak{}Reading}}\par\nobreak\noindent
The second way: grant every text, but not Paul’s reading on faith, nor Hebrews’, nor John’s. Faith is sufficient and works are not the ground, but faith is not shown to be necessary.

{\small\begin{verbatim}
def noNecessityWitnessReading : Valuation Claim :=
  fun a =>
    match a with
    | Claim.faithIsNecessary => False
    | Claim.justificationByFaithAlone => False
    | x => True
\end{verbatim}}
\par\medskip\noindent\textbf{\texttt{necessity\_\allowbreak{}witnesses\_\allowbreak{}jointly\_\allowbreak{}load\_\allowbreak{}bearing}}\par\nobreak\noindent
\textbf{To deny that faith is necessary, an opponent must reject three readings} — Paul’s on faith, Hebrews’ and John’s. Trent does not: it reads Paul’s “justified by faith” as faith the foundation of justification, without which it is impossible to please God. That part of faith alone is common ground.

{\small\begin{verbatim}
theorem necessity_witnesses_jointly_load_bearing : ¬Establishes
    withoutTheNecessityWitnesses
-- axioms: propext, Quot.sound
\end{verbatim}}
\par\medskip\noindent\textbf{\texttt{no\allowbreak{}Sufficiency\allowbreak{}Witness\allowbreak{}Reading}}\par\nobreak\noindent
The third way: grant every text, but not Paul’s reading on faith, nor John’s, Peter’s or Jesus’. Faith is necessary and works are not the ground, but faith is not shown to suffice.

{\small\begin{verbatim}
def noSufficiencyWitnessReading : Valuation Claim :=
  fun a =>
    match a with
    | Claim.faithIsSufficient => False
    | Claim.justificationByFaithAlone => False
    | x => True
\end{verbatim}}
\par\medskip\noindent\textbf{\texttt{sufficiency\_\allowbreak{}witnesses\_\allowbreak{}jointly\_\allowbreak{}load\_\allowbreak{}bearing}}\par\nobreak\noindent
\textbf{To deny that faith suffices, an opponent must reject four readings} — Paul’s on faith, John’s, Peter’s and Jesus’. It is the widest of the three, since four corpora say it. Trent’s canon 9 denies that faith suffices, but Trent’s position as encoded here (\texttt{tridentine}) takes the first way, at its reading of James 2:24.

{\small\begin{verbatim}
theorem sufficiency_witnesses_jointly_load_bearing : ¬Establishes
    withoutTheSufficiencyWitnesses
-- axioms: propext, Quot.sound
\end{verbatim}}
\par\medskip\noindent\textbf{\texttt{canonical\allowbreak{}With\allowbreak{}James\_\allowbreak{}burden}}\par\nobreak\noindent
\textbf{The opponent’s burden against the canonical witness with James.} The three ways of overturning faith alone remain, and a fourth is added for the Reformed formula: rejecting James’s reading alone. That set has one reading in it — James is the formula’s single point of failure, as he is its only witness that works are the fruit of faith.

{\small\begin{verbatim}
theorem canonicalWithJames_burden : OpponentsBurden canonicalWithJames
    (corpora ++ [jamesLine]) [p Claim.jude20_21] [toFaithAlone, toNeverAlone]
    [[Reading.paul, Reading.peter, Reading.jesus], [Reading.paulOnFaith,
    Reading.hebrews, Reading.john], [Reading.paulOnFaith, Reading.john,
    Reading.peter, Reading.jesus], [Reading.james]]
-- axioms: propext, Classical.choice, Quot.sound
\end{verbatim}}
\paragraph{What an opponent must reject.} \emph{The canonical witness, with James} concludes that justification is by faith alone, by a faith never alone. That conclusion is overturned exactly when every reading in one of these sets is rejected. Each set is minimal, and no other rejection overturns it:
\begin{enumerate}
\item \emph{Paul, on works (Romans 3:28, Ephesians 2:8–9, Titus 3:5)} (disputed); \emph{Peter (Acts 10 and 15)} (disputed); \emph{Jesus (Luke 7 and 18)} (disputed)
\item \emph{Paul, on faith (Romans 4:4–5, Galatians 2:16, 3:11)} (disputed); \emph{Hebrews} (consensus); \emph{John} (disputed)
\item \emph{Paul, on faith (Romans 4:4–5, Galatians 2:16, 3:11)} (disputed); \emph{John} (disputed); \emph{Peter (Acts 10 and 15)} (disputed); \emph{Jesus (Luke 7 and 18)} (disputed)
\item \emph{James 2} (disputed)
\end{enumerate}
\paragraph*{James}
\leavevmode\par\nobreak
\par\medskip\noindent\textbf{\texttt{james\allowbreak{}Undemonstrated\allowbreak{}Reading}}\par\nobreak\noindent
James read without the demonstrative sense of “justified”: works are not shown to be the fruit of faith, and the Reformed formula does not follow.

{\small\begin{verbatim}
def jamesUndemonstratedReading : Valuation Claim :=
  fun a =>
    match a with
    | Claim.jamesJustifiesDemonstratively => False
    | Claim.worksAreFruitOfFaith => False
    | Claim.faithAloneNeverAlone => False
    | x => True
\end{verbatim}}
\par\medskip\noindent\textbf{\texttt{james\_\allowbreak{}demonstrative\_\allowbreak{}is\_\allowbreak{}load\_\allowbreak{}bearing}}\par\nobreak\noindent
\textbf{The sense of “justified” in James is load-bearing} for the formula. Take the demonstrative reading away, keeping every text of James and the distinction between justification and sanctification, and “faith never alone” no longer follows.

{\small\begin{verbatim}
theorem james_demonstrative_is_load_bearing : ¬Establishes
    withJamesUndemonstrated
-- axioms: propext, Quot.sound
\end{verbatim}}
\paragraph*{Paul on faith, beside James}
\leavevmode\par\nobreak
Paul’s message about faith, stated on its own: faith is necessary and it is sufficient. Does James contradict it? Not on its texts: none of them says either. It turns on one reading — of whom James says “can that faith save him?” (2:14).

\par\medskip\noindent\textbf{\texttt{paul\_\allowbreak{}on\_\allowbreak{}faith\_\allowbreak{}establishes}}\par\nobreak\noindent
\textbf{Paul: faith is necessary and sufficient.} From Romans 4:4–5 and Galatians 2:16 and 3:11 alone, on Paul’s reading of them.

{\small\begin{verbatim}
theorem paul_on_faith_establishes : Establishes paulOnFaith
-- axioms: propext, Quot.sound
\end{verbatim}}
\par\medskip\noindent\textbf{\texttt{paul\allowbreak{}On\allowbreak{}Faith\_\allowbreak{}is\_\allowbreak{}satisfiable}}\par\nobreak\noindent
Paul’s claim about faith has a model.

{\small\begin{verbatim}
theorem paulOnFaith_is_satisfiable : Satisfiable paulOnFaith.premises
-- axioms: propext, Quot.sound
\end{verbatim}}
\par\medskip\noindent\textbf{\texttt{paul\allowbreak{}With\allowbreak{}James\allowbreak{}Plainly\allowbreak{}Reading}}\par\nobreak\noindent
James’s world read plainly, beside Paul’s: every text of James true, the faith James calls dead is faith without works, justification and sanctification distinct — and “justified” in 2:21–24 given its ordinary sense, not the demonstrative one. Faith is necessary and sufficient.

{\small\begin{verbatim}
def paulWithJamesPlainlyReading : Valuation Claim :=
  fun a =>
    match a with
    | Claim.jamesJustifiesDemonstratively => False
    | Claim.worksAreFruitOfFaith => False
    | Claim.faithAloneNeverAlone => False
    | x => True
\end{verbatim}}
\par\medskip\noindent\textbf{\texttt{paul\_\allowbreak{}on\_\allowbreak{}faith\_\allowbreak{}stands\_\allowbreak{}with\_\allowbreak{}james}}\par\nobreak\noindent
\textbf{Paul’s claim about faith stands with James.} Every text of James, the reading of 2:14–17 as about faith without works, and the distinction between justification and sanctification can all be held with faith necessary and sufficient — and without the demonstrative sense of “justified” on which James and Trent part. The harmony does not rest on that contested sense.

{\small\begin{verbatim}
theorem paul_on_faith_stands_with_james : Satisfiable (paulOnFaith.premises ++
    [paulOnFaith.conclusion] ++ jamesWithoutDemonstrative.premises)
-- axioms: propext, Quot.sound
\end{verbatim}}
\par\medskip\noindent\textbf{\texttt{james\allowbreak{}Against\allowbreak{}Sufficiency\allowbreak{}Reading}}\par\nobreak\noindent
The world of James read against sufficiency: every text of James true, and faith, as such, not enough.

{\small\begin{verbatim}
def jamesAgainstSufficiencyReading : Valuation Claim :=
  fun a =>
    match a with
    | Claim.faithIsSufficient => False
    | Claim.justificationByFaithAlone => False
    | Claim.faithAloneNeverAlone => False
    | x => True
\end{verbatim}}
\par\medskip\noindent\textbf{\texttt{james\allowbreak{}Read\allowbreak{}Against\allowbreak{}Sufficiency\_\allowbreak{}is\_\allowbreak{}satisfiable}}\par\nobreak\noindent
James read against sufficiency has a model: the rival is coherent.

{\small\begin{verbatim}
theorem jamesReadAgainstSufficiency_is_satisfiable : Satisfiable
    jamesReadAgainstSufficiency.premises
-- axioms: propext, Quot.sound
\end{verbatim}}
\par\medskip\noindent\textbf{\texttt{james\allowbreak{}Read\allowbreak{}Against\allowbreak{}Sufficiency\_\allowbreak{}establishes}}\par\nobreak\noindent
The rival reading establishes what it says: faith does not suffice.

{\small\begin{verbatim}
theorem jamesReadAgainstSufficiency_establishes : Establishes
    jamesReadAgainstSufficiency
-- axioms: propext, Quot.sound
\end{verbatim}}
\par\medskip\noindent\textbf{\texttt{james\_\allowbreak{}against\_\allowbreak{}sufficiency\_\allowbreak{}contradicts\_\allowbreak{}paul}}\par\nobreak\noindent
\textbf{What the harmony rests on.} Read “that faith” in James 2:14 and “not by faith alone” in 2:24 of faith as such, as Trent does, and James and Paul’s claim about faith cannot both be held. So the harmony is a reading, not a given: it rests on the faith James says cannot save being the faith 2:14–17 describes — faith without works — which is \texttt{james2\allowbreak{}Targets\allowbreak{}Dead\allowbreak{}Faith}, cited to Moo and Johnson.

{\small\begin{verbatim}
theorem james_against_sufficiency_contradicts_paul : ¬Satisfiable
    (paulOnFaith.premises ++ jamesReadAgainstSufficiency.premises)
-- axioms: propext, Classical.choice, Quot.sound
\end{verbatim}}
\paragraph*{Why the canonical witness stands against Trent}
\leavevmode\par\nobreak
Both positions hold every text. They part at one reading: the sense of “justified” in James 2:21–24.

\par\medskip\noindent\textbf{\texttt{why\allowbreak{}The\allowbreak{}Canonical\allowbreak{}Witness\allowbreak{}Stands\allowbreak{}Against\allowbreak{}Trent}}\par\nobreak\noindent
\textbf{Why the canonical witness, with James, stands against Trent.} The crux is the demonstrative sense of “justified” in James: it breaks Trent at exactly its reading of James 2:24, and the Reformed formula needs it. Every other premise of Trent — every text — the canonical witness holds too.

\paragraph{Why \emph{The canonical witness, with James} stands against \emph{Trent, reading James 2:24 as justification increased by works}.}
\begin{itemize}
\item \textbf{The crux:} \(\pvar{TestimonyArguments.CanonicalWitness}{28}\), a premise of \emph{The canonical witness, with James}.
\item \textbf{What it does:} its conclusion does not follow without it.
\item \textbf{Where the rival breaks:} \(\lnot \pvar{TestimonyArguments.CanonicalWitness}{28}\) cannot be held together with \(\pvar{TestimonyArguments.CanonicalWitness}{28}\); each is needed for the break, and without the crux they stand.
\item \textbf{What this rests on:}
\begin{itemize}
\item \(\pvar{TestimonyArguments.CanonicalWitness}{28}\) ‘Justified’ in James 2:21–24 means shown to be righteous, not made righteous --- \emph{disputed}: \cite[III.xvii.11]{calvin-institutes-1960}; Jas 2:21-23
\end{itemize}
\end{itemize}
\paragraph*{How the burdens are built}
\leavevmode\par\nobreak
Each is built with \texttt{Argument\allowbreak{}Package.rejecting}, the operation the opponent’s burden quantifies over (\texttt{Testimony.Logic.Burden}), and its result is proved from that burden in \texttt{Testimony\allowbreak{}Arguments.Canonical\allowbreak{}Witness.Results}. So the named burdens are instances of the computed one, and cannot drift from it.

\par\medskip\noindent\textbf{\texttt{all\allowbreak{}Texts}}\par\nobreak\noindent
Every text the argument cites, from every corpus and from James. The rival holds all of them.

\begin{derivation}
  \item[(1)] \(\pvar{TestimonyArguments.CanonicalWitness}{1}\)\gloss{Romans 3:28: a person is justified by faith apart from works of the law}
  \item[(2)] \(\pvar{TestimonyArguments.CanonicalWitness}{5}\)\gloss{Ephesians 2:8–9: saved by grace through faith, not of works}
  \item[(3)] \(\pvar{TestimonyArguments.CanonicalWitness}{6}\)\gloss{Titus 3:5: saved not by works done in righteousness but by his mercy}
  \breakrow{(4)} \(\pvar{TestimonyArguments.CanonicalWitness}{2}\)\gloss{Romans 4:4–5: to the one who does not work but believes, faith is counted as righteousness}
  \breakrow{(5)} \(\pvar{TestimonyArguments.CanonicalWitness}{3}\)\gloss{Galatians 2:16: not justified by works of the law but through faith in Christ}
  \breakrow{(6)} \(\pvar{TestimonyArguments.CanonicalWitness}{4}\)\gloss{Galatians 3:11: no one is justified before God by the law; the righteous shall live by faith}
  \breakrow{(7)} \(\pvar{TestimonyArguments.CanonicalWitness}{7}\)\gloss{Hebrews 10:38–39: the righteous one shall live by faith}
  \breakrow{(8)} \(\pvar{TestimonyArguments.CanonicalWitness}{8}\)\gloss{Hebrews 11:6: without faith it is impossible to please God}
  \breakrow{(9)} \(\pvar{TestimonyArguments.CanonicalWitness}{9}\)\gloss{John 3:16–18: whoever believes has eternal life; whoever does not believe is condemned already}
  \breakrow{(10)} \(\pvar{TestimonyArguments.CanonicalWitness}{10}\)\gloss{John 5:24: whoever hears and believes has eternal life and does not come into judgement}
  \breakrow{(11)} \(\pvar{TestimonyArguments.CanonicalWitness}{11}\)\gloss{John 6:28–29: the work of God is that you believe in him whom he has sent}
  \breakrow{(12)} \(\pvar{TestimonyArguments.CanonicalWitness}{12}\)\gloss{Acts 10:43: everyone who believes in him receives forgiveness of sins}
  \breakrow{(13)} \(\pvar{TestimonyArguments.CanonicalWitness}{13}\)\gloss{Acts 15:10–11: no yoke on the disciples; we are saved through the grace of the Lord Jesus}
  \breakrow{(14)} \(\pvar{TestimonyArguments.CanonicalWitness}{14}\)\gloss{Luke 7:50: your faith has saved you}
  \breakrow{(15)} \(\pvar{TestimonyArguments.CanonicalWitness}{15}\)\gloss{Luke 18:9–14: the tax collector, not the Pharisee who listed his works, went home justified}
  \breakrow{(16)} \(\pvar{TestimonyArguments.CanonicalWitness}{16}\)\gloss{James 2:14–17: faith by itself, if it does not have works, is dead}
  \breakrow{(17)} \(\pvar{TestimonyArguments.CanonicalWitness}{17}\)\gloss{James 2:19: even the demons believe, and shudder}
  \breakrow{(18)} \(\pvar{TestimonyArguments.CanonicalWitness}{18}\)\gloss{James 2:21–23: Abraham’s faith was completed by his works, and Genesis 15:6 was fulfilled}
  \item[(19)] \(\pvar{TestimonyArguments.CanonicalWitness}{19}\)\gloss{James 2:24: a person is justified by works and not by faith alone}
  \item[(20)] \(\pvar{TestimonyArguments.CanonicalWitness}{20}\)\gloss{Jude 20–21: keep yourselves in the love of God, waiting for the mercy of our Lord Jesus Christ}
\end{derivation}
\par\medskip\noindent\textbf{\texttt{tridentine}} --- Trent, reading James 2:24 as justification increased by works\par\nobreak\noindent
\textbf{Trent.} It grants every text the canonical witness reads — its Decree on Justification quotes most of them — and reads James 2:24 as the increase of a justification already received, by works done in grace (Session VI, ch. 10). So works are part of the ground, and faith alone is denied.

\begin{derivation}
  \item[(1)] \(\pvar{TestimonyArguments.CanonicalWitness}{1}\)\gloss{Romans 3:28: a person is justified by faith apart from works of the law}
  \item[(2)] \(\pvar{TestimonyArguments.CanonicalWitness}{5}\)\gloss{Ephesians 2:8–9: saved by grace through faith, not of works}
  \item[(3)] \(\pvar{TestimonyArguments.CanonicalWitness}{6}\)\gloss{Titus 3:5: saved not by works done in righteousness but by his mercy}
  \breakrow{(4)} \(\pvar{TestimonyArguments.CanonicalWitness}{2}\)\gloss{Romans 4:4–5: to the one who does not work but believes, faith is counted as righteousness}
  \breakrow{(5)} \(\pvar{TestimonyArguments.CanonicalWitness}{3}\)\gloss{Galatians 2:16: not justified by works of the law but through faith in Christ}
  \breakrow{(6)} \(\pvar{TestimonyArguments.CanonicalWitness}{4}\)\gloss{Galatians 3:11: no one is justified before God by the law; the righteous shall live by faith}
  \breakrow{(7)} \(\pvar{TestimonyArguments.CanonicalWitness}{7}\)\gloss{Hebrews 10:38–39: the righteous one shall live by faith}
  \breakrow{(8)} \(\pvar{TestimonyArguments.CanonicalWitness}{8}\)\gloss{Hebrews 11:6: without faith it is impossible to please God}
  \breakrow{(9)} \(\pvar{TestimonyArguments.CanonicalWitness}{9}\)\gloss{John 3:16–18: whoever believes has eternal life; whoever does not believe is condemned already}
  \breakrow{(10)} \(\pvar{TestimonyArguments.CanonicalWitness}{10}\)\gloss{John 5:24: whoever hears and believes has eternal life and does not come into judgement}
  \breakrow{(11)} \(\pvar{TestimonyArguments.CanonicalWitness}{11}\)\gloss{John 6:28–29: the work of God is that you believe in him whom he has sent}
  \breakrow{(12)} \(\pvar{TestimonyArguments.CanonicalWitness}{12}\)\gloss{Acts 10:43: everyone who believes in him receives forgiveness of sins}
  \breakrow{(13)} \(\pvar{TestimonyArguments.CanonicalWitness}{13}\)\gloss{Acts 15:10–11: no yoke on the disciples; we are saved through the grace of the Lord Jesus}
  \breakrow{(14)} \(\pvar{TestimonyArguments.CanonicalWitness}{14}\)\gloss{Luke 7:50: your faith has saved you}
  \breakrow{(15)} \(\pvar{TestimonyArguments.CanonicalWitness}{15}\)\gloss{Luke 18:9–14: the tax collector, not the Pharisee who listed his works, went home justified}
  \breakrow{(16)} \(\pvar{TestimonyArguments.CanonicalWitness}{16}\)\gloss{James 2:14–17: faith by itself, if it does not have works, is dead}
  \breakrow{(17)} \(\pvar{TestimonyArguments.CanonicalWitness}{17}\)\gloss{James 2:19: even the demons believe, and shudder}
  \breakrow{(18)} \(\pvar{TestimonyArguments.CanonicalWitness}{18}\)\gloss{James 2:21–23: Abraham’s faith was completed by his works, and Genesis 15:6 was fulfilled}
  \breakrow{(19)} \(\pvar{TestimonyArguments.CanonicalWitness}{19}\)\gloss{James 2:24: a person is justified by works and not by faith alone}
  \breakrow{(20)} \(\pvar{TestimonyArguments.CanonicalWitness}{20}\)\gloss{Jude 20–21: keep yourselves in the love of God, waiting for the mercy of our Lord Jesus Christ}
  \breakrow{(21)} \(\lnot \pvar{TestimonyArguments.CanonicalWitness}{28}\)\gloss{\emph{denied:} ‘Justified’ in James 2:21–24 means shown to be righteous (denied by Trent, which reads 2:24 as an increase of justification by works done in grace)}
  \breakrow{(22)} \((\pvar{TestimonyArguments.CanonicalWitness}{19} \land \lnot \pvar{TestimonyArguments.CanonicalWitness}{28}) \rightarrow \lnot \pvar{TestimonyArguments.CanonicalWitness}{23}\)\gloss{\emph{would deny:} Works are not the ground of justification}
  \item[(23)] \(\lnot \pvar{TestimonyArguments.CanonicalWitness}{23} \rightarrow \lnot \pvar{TestimonyArguments.CanonicalWitness}{24}\)
  \item[\(\vdash\)] \(\lnot \pvar{TestimonyArguments.CanonicalWitness}{24}\)\gloss{\emph{at issue, the denial of:} Justification is by faith alone}
  \keepwithnext
\end{derivation}
\emph{Grounded in scripture alone:} Romans 3:28: a person is justified by faith apart from works of the law, Ephesians 2:8–9: saved by grace through faith, not of works, Titus 3:5: saved not by works done in righteousness but by his mercy, Romans 4:4–5: to the one who does not work but believes, faith is counted as righteousness, Galatians 2:16: not justified by works of the law but through faith in Christ, Galatians 3:11: no one is justified before God by the law; the righteous shall live by faith, Hebrews 10:38–39: the righteous one shall live by faith, Hebrews 11:6: without faith it is impossible to please God, John 3:16–18: whoever believes has eternal life; whoever does not believe is condemned already, John 5:24: whoever hears and believes has eternal life and does not come into judgement, John 6:28–29: the work of God is that you believe in him whom he has sent, Acts 10:43: everyone who believes in him receives forgiveness of sins, Acts 15:10–11: no yoke on the disciples; we are saved through the grace of the Lord Jesus, Luke 7:50: your faith has saved you, Luke 18:9–14: the tax collector, not the Pharisee who listed his works, went home justified, James 2:14–17: faith by itself, if it does not have works, is dead, James 2:19: even the demons believe, and shudder, James 2:21–23: Abraham’s faith was completed by his works, and Genesis 15:6 was fulfilled, James 2:24: a person is justified by works and not by faith alone, Jude 20–21: keep yourselves in the love of God, waiting for the mercy of our Lord Jesus Christ.
\par\medskip\noindent\textbf{\texttt{canonical\allowbreak{}Case}} --- The canonical witness to faith alone\par\nobreak\noindent
\textbf{The canonical witness}: every corpus, each on its own reading, and the three parts of faith alone put together.

\begin{derivation}
  \item[(1)] \(\pvar{TestimonyArguments.CanonicalWitness}{1}\)\gloss{Romans 3:28: a person is justified by faith apart from works of the law}
  \item[(2)] \(\pvar{TestimonyArguments.CanonicalWitness}{5}\)\gloss{Ephesians 2:8–9: saved by grace through faith, not of works}
  \item[(3)] \(\pvar{TestimonyArguments.CanonicalWitness}{6}\)\gloss{Titus 3:5: saved not by works done in righteousness but by his mercy}
  \breakrow{(4)} \(\pvar{TestimonyArguments.CanonicalWitness}{2}\)\gloss{Romans 4:4–5: to the one who does not work but believes, faith is counted as righteousness}
  \breakrow{(5)} \(\pvar{TestimonyArguments.CanonicalWitness}{3}\)\gloss{Galatians 2:16: not justified by works of the law but through faith in Christ}
  \breakrow{(6)} \(\pvar{TestimonyArguments.CanonicalWitness}{4}\)\gloss{Galatians 3:11: no one is justified before God by the law; the righteous shall live by faith}
  \breakrow{(7)} \(\pvar{TestimonyArguments.CanonicalWitness}{7}\)\gloss{Hebrews 10:38–39: the righteous one shall live by faith}
  \breakrow{(8)} \(\pvar{TestimonyArguments.CanonicalWitness}{8}\)\gloss{Hebrews 11:6: without faith it is impossible to please God}
  \breakrow{(9)} \(\pvar{TestimonyArguments.CanonicalWitness}{9}\)\gloss{John 3:16–18: whoever believes has eternal life; whoever does not believe is condemned already}
  \breakrow{(10)} \(\pvar{TestimonyArguments.CanonicalWitness}{10}\)\gloss{John 5:24: whoever hears and believes has eternal life and does not come into judgement}
  \breakrow{(11)} \(\pvar{TestimonyArguments.CanonicalWitness}{11}\)\gloss{John 6:28–29: the work of God is that you believe in him whom he has sent}
  \breakrow{(12)} \(\pvar{TestimonyArguments.CanonicalWitness}{12}\)\gloss{Acts 10:43: everyone who believes in him receives forgiveness of sins}
  \breakrow{(13)} \(\pvar{TestimonyArguments.CanonicalWitness}{13}\)\gloss{Acts 15:10–11: no yoke on the disciples; we are saved through the grace of the Lord Jesus}
  \breakrow{(14)} \(\pvar{TestimonyArguments.CanonicalWitness}{14}\)\gloss{Luke 7:50: your faith has saved you}
  \breakrow{(15)} \(\pvar{TestimonyArguments.CanonicalWitness}{15}\)\gloss{Luke 18:9–14: the tax collector, not the Pharisee who listed his works, went home justified}
  \breakrow{(16)} \((\pvar{TestimonyArguments.CanonicalWitness}{1} \land \pvar{TestimonyArguments.CanonicalWitness}{5} \land \pvar{TestimonyArguments.CanonicalWitness}{6}) \rightarrow \pvar{TestimonyArguments.CanonicalWitness}{23}\)\gloss{\emph{would give:} Works are not the ground of justification}
  \breakrow{(17)} \((\pvar{TestimonyArguments.CanonicalWitness}{2} \land \pvar{TestimonyArguments.CanonicalWitness}{3} \land \pvar{TestimonyArguments.CanonicalWitness}{4}) \rightarrow \pvar{TestimonyArguments.CanonicalWitness}{21} \land \pvar{TestimonyArguments.CanonicalWitness}{22}\)\gloss{\emph{would give:} \(\pvar{TestimonyArguments.CanonicalWitness}{21}\) Faith is necessary: no one is saved without it; \(\pvar{TestimonyArguments.CanonicalWitness}{22}\) Faith is sufficient: whoever believes is saved}
  \breakrow{(18)} \((\pvar{TestimonyArguments.CanonicalWitness}{7} \land \pvar{TestimonyArguments.CanonicalWitness}{8}) \rightarrow \pvar{TestimonyArguments.CanonicalWitness}{21}\)\gloss{\emph{would give:} Faith is necessary: no one is saved without it}
  \breakrow{(19)} \((\pvar{TestimonyArguments.CanonicalWitness}{9} \land \pvar{TestimonyArguments.CanonicalWitness}{10} \land \pvar{TestimonyArguments.CanonicalWitness}{11}) \rightarrow \pvar{TestimonyArguments.CanonicalWitness}{22} \land \pvar{TestimonyArguments.CanonicalWitness}{21}\)\gloss{\emph{would give:} \(\pvar{TestimonyArguments.CanonicalWitness}{22}\) Faith is sufficient: whoever believes is saved; \(\pvar{TestimonyArguments.CanonicalWitness}{21}\) Faith is necessary: no one is saved without it}
  \breakrow{(20)} \((\pvar{TestimonyArguments.CanonicalWitness}{12} \land \pvar{TestimonyArguments.CanonicalWitness}{13}) \rightarrow \pvar{TestimonyArguments.CanonicalWitness}{22} \land \pvar{TestimonyArguments.CanonicalWitness}{23}\)\gloss{\emph{would give:} \(\pvar{TestimonyArguments.CanonicalWitness}{22}\) Faith is sufficient: whoever believes is saved; \(\pvar{TestimonyArguments.CanonicalWitness}{23}\) Works are not the ground of justification}
  \breakrow{(21)} \((\pvar{TestimonyArguments.CanonicalWitness}{14} \land \pvar{TestimonyArguments.CanonicalWitness}{15}) \rightarrow \pvar{TestimonyArguments.CanonicalWitness}{22} \land \pvar{TestimonyArguments.CanonicalWitness}{23}\)\gloss{\emph{would give:} \(\pvar{TestimonyArguments.CanonicalWitness}{22}\) Faith is sufficient: whoever believes is saved; \(\pvar{TestimonyArguments.CanonicalWitness}{23}\) Works are not the ground of justification}
  \item[(22)] \((\pvar{TestimonyArguments.CanonicalWitness}{21} \land \pvar{TestimonyArguments.CanonicalWitness}{22} \land \pvar{TestimonyArguments.CanonicalWitness}{23}) \rightarrow \pvar{TestimonyArguments.CanonicalWitness}{24}\)
  \item[\(\vdash\)] \(\pvar{TestimonyArguments.CanonicalWitness}{24}\)\gloss{\emph{at issue:} Justification is by faith alone}
  \keepwithnext
\end{derivation}
\emph{Grounded in scripture alone:} Romans 3:28: a person is justified by faith apart from works of the law, Ephesians 2:8–9: saved by grace through faith, not of works, Titus 3:5: saved not by works done in righteousness but by his mercy, Romans 4:4–5: to the one who does not work but believes, faith is counted as righteousness, Galatians 2:16: not justified by works of the law but through faith in Christ, Galatians 3:11: no one is justified before God by the law; the righteous shall live by faith, Hebrews 10:38–39: the righteous one shall live by faith, Hebrews 11:6: without faith it is impossible to please God, John 3:16–18: whoever believes has eternal life; whoever does not believe is condemned already, John 5:24: whoever hears and believes has eternal life and does not come into judgement, John 6:28–29: the work of God is that you believe in him whom he has sent, Acts 10:43: everyone who believes in him receives forgiveness of sins, Acts 15:10–11: no yoke on the disciples; we are saved through the grace of the Lord Jesus, Luke 7:50: your faith has saved you, Luke 18:9–14: the tax collector, not the Pharisee who listed his works, went home justified.
\par\medskip\noindent\textbf{\texttt{canonical\allowbreak{}With\allowbreak{}James}} --- The canonical witness, with James\par\nobreak\noindent
\textbf{The canonical witness with James}: James read with the distinction between justification and sanctification, and Jude as consistent with the rest. It concludes the Reformed formula — faith alone, but a faith never alone.

\begin{derivation}
  \item[(1)] \(\pvar{TestimonyArguments.CanonicalWitness}{1}\)\gloss{Romans 3:28: a person is justified by faith apart from works of the law}
  \item[(2)] \(\pvar{TestimonyArguments.CanonicalWitness}{5}\)\gloss{Ephesians 2:8–9: saved by grace through faith, not of works}
  \item[(3)] \(\pvar{TestimonyArguments.CanonicalWitness}{6}\)\gloss{Titus 3:5: saved not by works done in righteousness but by his mercy}
  \breakrow{(4)} \(\pvar{TestimonyArguments.CanonicalWitness}{2}\)\gloss{Romans 4:4–5: to the one who does not work but believes, faith is counted as righteousness}
  \breakrow{(5)} \(\pvar{TestimonyArguments.CanonicalWitness}{3}\)\gloss{Galatians 2:16: not justified by works of the law but through faith in Christ}
  \breakrow{(6)} \(\pvar{TestimonyArguments.CanonicalWitness}{4}\)\gloss{Galatians 3:11: no one is justified before God by the law; the righteous shall live by faith}
  \breakrow{(7)} \(\pvar{TestimonyArguments.CanonicalWitness}{7}\)\gloss{Hebrews 10:38–39: the righteous one shall live by faith}
  \breakrow{(8)} \(\pvar{TestimonyArguments.CanonicalWitness}{8}\)\gloss{Hebrews 11:6: without faith it is impossible to please God}
  \breakrow{(9)} \(\pvar{TestimonyArguments.CanonicalWitness}{9}\)\gloss{John 3:16–18: whoever believes has eternal life; whoever does not believe is condemned already}
  \breakrow{(10)} \(\pvar{TestimonyArguments.CanonicalWitness}{10}\)\gloss{John 5:24: whoever hears and believes has eternal life and does not come into judgement}
  \breakrow{(11)} \(\pvar{TestimonyArguments.CanonicalWitness}{11}\)\gloss{John 6:28–29: the work of God is that you believe in him whom he has sent}
  \breakrow{(12)} \(\pvar{TestimonyArguments.CanonicalWitness}{12}\)\gloss{Acts 10:43: everyone who believes in him receives forgiveness of sins}
  \breakrow{(13)} \(\pvar{TestimonyArguments.CanonicalWitness}{13}\)\gloss{Acts 15:10–11: no yoke on the disciples; we are saved through the grace of the Lord Jesus}
  \breakrow{(14)} \(\pvar{TestimonyArguments.CanonicalWitness}{14}\)\gloss{Luke 7:50: your faith has saved you}
  \breakrow{(15)} \(\pvar{TestimonyArguments.CanonicalWitness}{15}\)\gloss{Luke 18:9–14: the tax collector, not the Pharisee who listed his works, went home justified}
  \breakrow{(16)} \(\pvar{TestimonyArguments.CanonicalWitness}{16}\)\gloss{James 2:14–17: faith by itself, if it does not have works, is dead}
  \breakrow{(17)} \(\pvar{TestimonyArguments.CanonicalWitness}{17}\)\gloss{James 2:19: even the demons believe, and shudder}
  \breakrow{(18)} \(\pvar{TestimonyArguments.CanonicalWitness}{18}\)\gloss{James 2:21–23: Abraham’s faith was completed by his works, and Genesis 15:6 was fulfilled}
  \breakrow{(19)} \(\pvar{TestimonyArguments.CanonicalWitness}{19}\)\gloss{James 2:24: a person is justified by works and not by faith alone}
  \breakrow{(20)} \(\pvar{TestimonyArguments.CanonicalWitness}{27}\)\gloss{James 2:14–26 targets a barren faith — mere assent — not Paul’s doctrine}
  \breakrow{(21)} \(\pvar{TestimonyArguments.CanonicalWitness}{28}\)\gloss{‘Justified’ in James 2:21–24 means shown to be righteous (denied by Trent, which reads 2:24 as an increase of justification by works done in grace)}
  \breakrow{(22)} \(\pvar{TestimonyArguments.CanonicalWitness}{29}\)\gloss{Justification and sanctification are inseparable but distinct}
  \breakrow{(23)} \(\pvar{TestimonyArguments.CanonicalWitness}{20}\)\gloss{Jude 20–21: keep yourselves in the love of God, waiting for the mercy of our Lord Jesus Christ}
  \breakrow{(24)} \((\pvar{TestimonyArguments.CanonicalWitness}{1} \land \pvar{TestimonyArguments.CanonicalWitness}{5} \land \pvar{TestimonyArguments.CanonicalWitness}{6}) \rightarrow \pvar{TestimonyArguments.CanonicalWitness}{23}\)\gloss{\emph{would give:} Works are not the ground of justification}
  \breakrow{(25)} \((\pvar{TestimonyArguments.CanonicalWitness}{2} \land \pvar{TestimonyArguments.CanonicalWitness}{3} \land \pvar{TestimonyArguments.CanonicalWitness}{4}) \rightarrow \pvar{TestimonyArguments.CanonicalWitness}{21} \land \pvar{TestimonyArguments.CanonicalWitness}{22}\)\gloss{\emph{would give:} \(\pvar{TestimonyArguments.CanonicalWitness}{21}\) Faith is necessary: no one is saved without it; \(\pvar{TestimonyArguments.CanonicalWitness}{22}\) Faith is sufficient: whoever believes is saved}
  \breakrow{(26)} \((\pvar{TestimonyArguments.CanonicalWitness}{7} \land \pvar{TestimonyArguments.CanonicalWitness}{8}) \rightarrow \pvar{TestimonyArguments.CanonicalWitness}{21}\)\gloss{\emph{would give:} Faith is necessary: no one is saved without it}
  \breakrow{(27)} \((\pvar{TestimonyArguments.CanonicalWitness}{9} \land \pvar{TestimonyArguments.CanonicalWitness}{10} \land \pvar{TestimonyArguments.CanonicalWitness}{11}) \rightarrow \pvar{TestimonyArguments.CanonicalWitness}{22} \land \pvar{TestimonyArguments.CanonicalWitness}{21}\)\gloss{\emph{would give:} \(\pvar{TestimonyArguments.CanonicalWitness}{22}\) Faith is sufficient: whoever believes is saved; \(\pvar{TestimonyArguments.CanonicalWitness}{21}\) Faith is necessary: no one is saved without it}
  \breakrow{(28)} \((\pvar{TestimonyArguments.CanonicalWitness}{12} \land \pvar{TestimonyArguments.CanonicalWitness}{13}) \rightarrow \pvar{TestimonyArguments.CanonicalWitness}{22} \land \pvar{TestimonyArguments.CanonicalWitness}{23}\)\gloss{\emph{would give:} \(\pvar{TestimonyArguments.CanonicalWitness}{22}\) Faith is sufficient: whoever believes is saved; \(\pvar{TestimonyArguments.CanonicalWitness}{23}\) Works are not the ground of justification}
  \breakrow{(29)} \((\pvar{TestimonyArguments.CanonicalWitness}{14} \land \pvar{TestimonyArguments.CanonicalWitness}{15}) \rightarrow \pvar{TestimonyArguments.CanonicalWitness}{22} \land \pvar{TestimonyArguments.CanonicalWitness}{23}\)\gloss{\emph{would give:} \(\pvar{TestimonyArguments.CanonicalWitness}{22}\) Faith is sufficient: whoever believes is saved; \(\pvar{TestimonyArguments.CanonicalWitness}{23}\) Works are not the ground of justification}
  \breakrow{(30)} \((\pvar{TestimonyArguments.CanonicalWitness}{16} \land \pvar{TestimonyArguments.CanonicalWitness}{18} \land \pvar{TestimonyArguments.CanonicalWitness}{27} \land \pvar{TestimonyArguments.CanonicalWitness}{28} \land \pvar{TestimonyArguments.CanonicalWitness}{29}) \rightarrow \pvar{TestimonyArguments.CanonicalWitness}{25}\)\gloss{\emph{would give:} Good works are the fruit and evidence of a true and lively faith}
  \breakrow{(31)} \((\pvar{TestimonyArguments.CanonicalWitness}{21} \land \pvar{TestimonyArguments.CanonicalWitness}{22} \land \pvar{TestimonyArguments.CanonicalWitness}{23}) \rightarrow \pvar{TestimonyArguments.CanonicalWitness}{24}\)\gloss{\emph{would give:} Justification is by faith alone}
  \item[(32)] \((\pvar{TestimonyArguments.CanonicalWitness}{24} \land \pvar{TestimonyArguments.CanonicalWitness}{25}) \rightarrow \pvar{TestimonyArguments.CanonicalWitness}{26}\)
  \item[\(\vdash\)] \(\pvar{TestimonyArguments.CanonicalWitness}{26}\)\gloss{\emph{at issue:} Justification is by faith alone, by a faith that is never alone}
  \keepwithnext
\end{derivation}
\emph{Grounded in scripture alone:} Romans 3:28: a person is justified by faith apart from works of the law, Ephesians 2:8–9: saved by grace through faith, not of works, Titus 3:5: saved not by works done in righteousness but by his mercy, Romans 4:4–5: to the one who does not work but believes, faith is counted as righteousness, Galatians 2:16: not justified by works of the law but through faith in Christ, Galatians 3:11: no one is justified before God by the law; the righteous shall live by faith, Hebrews 10:38–39: the righteous one shall live by faith, Hebrews 11:6: without faith it is impossible to please God, John 3:16–18: whoever believes has eternal life; whoever does not believe is condemned already, John 5:24: whoever hears and believes has eternal life and does not come into judgement, John 6:28–29: the work of God is that you believe in him whom he has sent, Acts 10:43: everyone who believes in him receives forgiveness of sins, Acts 15:10–11: no yoke on the disciples; we are saved through the grace of the Lord Jesus, Luke 7:50: your faith has saved you, Luke 18:9–14: the tax collector, not the Pharisee who listed his works, went home justified, James 2:14–17: faith by itself, if it does not have works, is dead, James 2:19: even the demons believe, and shudder, James 2:21–23: Abraham’s faith was completed by his works, and Genesis 15:6 was fulfilled, James 2:24: a person is justified by works and not by faith alone, Jude 20–21: keep yourselves in the love of God, waiting for the mercy of our Lord Jesus Christ.
\paragraph*{One corpus’s reading rejected}
\leavevmode\par\nobreak
\par\medskip\noindent\textbf{\texttt{without\allowbreak{}Paul}} --- The canonical witness, without Paul’s readings\par\nobreak\noindent
Without Paul: both his readings, on works and on faith.

\begin{derivation}
  \item[(1)] \(\pvar{TestimonyArguments.CanonicalWitness}{1}\)\gloss{Romans 3:28: a person is justified by faith apart from works of the law}
  \item[(2)] \(\pvar{TestimonyArguments.CanonicalWitness}{5}\)\gloss{Ephesians 2:8–9: saved by grace through faith, not of works}
  \item[(3)] \(\pvar{TestimonyArguments.CanonicalWitness}{6}\)\gloss{Titus 3:5: saved not by works done in righteousness but by his mercy}
  \breakrow{(4)} \(\pvar{TestimonyArguments.CanonicalWitness}{2}\)\gloss{Romans 4:4–5: to the one who does not work but believes, faith is counted as righteousness}
  \breakrow{(5)} \(\pvar{TestimonyArguments.CanonicalWitness}{3}\)\gloss{Galatians 2:16: not justified by works of the law but through faith in Christ}
  \breakrow{(6)} \(\pvar{TestimonyArguments.CanonicalWitness}{4}\)\gloss{Galatians 3:11: no one is justified before God by the law; the righteous shall live by faith}
  \breakrow{(7)} \(\pvar{TestimonyArguments.CanonicalWitness}{7}\)\gloss{Hebrews 10:38–39: the righteous one shall live by faith}
  \breakrow{(8)} \(\pvar{TestimonyArguments.CanonicalWitness}{8}\)\gloss{Hebrews 11:6: without faith it is impossible to please God}
  \breakrow{(9)} \(\pvar{TestimonyArguments.CanonicalWitness}{9}\)\gloss{John 3:16–18: whoever believes has eternal life; whoever does not believe is condemned already}
  \breakrow{(10)} \(\pvar{TestimonyArguments.CanonicalWitness}{10}\)\gloss{John 5:24: whoever hears and believes has eternal life and does not come into judgement}
  \breakrow{(11)} \(\pvar{TestimonyArguments.CanonicalWitness}{11}\)\gloss{John 6:28–29: the work of God is that you believe in him whom he has sent}
  \breakrow{(12)} \(\pvar{TestimonyArguments.CanonicalWitness}{12}\)\gloss{Acts 10:43: everyone who believes in him receives forgiveness of sins}
  \breakrow{(13)} \(\pvar{TestimonyArguments.CanonicalWitness}{13}\)\gloss{Acts 15:10–11: no yoke on the disciples; we are saved through the grace of the Lord Jesus}
  \breakrow{(14)} \(\pvar{TestimonyArguments.CanonicalWitness}{14}\)\gloss{Luke 7:50: your faith has saved you}
  \breakrow{(15)} \(\pvar{TestimonyArguments.CanonicalWitness}{15}\)\gloss{Luke 18:9–14: the tax collector, not the Pharisee who listed his works, went home justified}
  \breakrow{(16)} \((\pvar{TestimonyArguments.CanonicalWitness}{7} \land \pvar{TestimonyArguments.CanonicalWitness}{8}) \rightarrow \pvar{TestimonyArguments.CanonicalWitness}{21}\)\gloss{\emph{would give:} Faith is necessary: no one is saved without it}
  \breakrow{(17)} \((\pvar{TestimonyArguments.CanonicalWitness}{9} \land \pvar{TestimonyArguments.CanonicalWitness}{10} \land \pvar{TestimonyArguments.CanonicalWitness}{11}) \rightarrow \pvar{TestimonyArguments.CanonicalWitness}{22} \land \pvar{TestimonyArguments.CanonicalWitness}{21}\)\gloss{\emph{would give:} \(\pvar{TestimonyArguments.CanonicalWitness}{22}\) Faith is sufficient: whoever believes is saved; \(\pvar{TestimonyArguments.CanonicalWitness}{21}\) Faith is necessary: no one is saved without it}
  \breakrow{(18)} \((\pvar{TestimonyArguments.CanonicalWitness}{12} \land \pvar{TestimonyArguments.CanonicalWitness}{13}) \rightarrow \pvar{TestimonyArguments.CanonicalWitness}{22} \land \pvar{TestimonyArguments.CanonicalWitness}{23}\)\gloss{\emph{would give:} \(\pvar{TestimonyArguments.CanonicalWitness}{22}\) Faith is sufficient: whoever believes is saved; \(\pvar{TestimonyArguments.CanonicalWitness}{23}\) Works are not the ground of justification}
  \breakrow{(19)} \((\pvar{TestimonyArguments.CanonicalWitness}{14} \land \pvar{TestimonyArguments.CanonicalWitness}{15}) \rightarrow \pvar{TestimonyArguments.CanonicalWitness}{22} \land \pvar{TestimonyArguments.CanonicalWitness}{23}\)\gloss{\emph{would give:} \(\pvar{TestimonyArguments.CanonicalWitness}{22}\) Faith is sufficient: whoever believes is saved; \(\pvar{TestimonyArguments.CanonicalWitness}{23}\) Works are not the ground of justification}
  \item[(20)] \((\pvar{TestimonyArguments.CanonicalWitness}{21} \land \pvar{TestimonyArguments.CanonicalWitness}{22} \land \pvar{TestimonyArguments.CanonicalWitness}{23}) \rightarrow \pvar{TestimonyArguments.CanonicalWitness}{24}\)
  \item[\(\vdash\)] \(\pvar{TestimonyArguments.CanonicalWitness}{24}\)\gloss{\emph{at issue:} Justification is by faith alone}
  \keepwithnext
\end{derivation}
\emph{Grounded in scripture alone:} Romans 3:28: a person is justified by faith apart from works of the law, Ephesians 2:8–9: saved by grace through faith, not of works, Titus 3:5: saved not by works done in righteousness but by his mercy, Romans 4:4–5: to the one who does not work but believes, faith is counted as righteousness, Galatians 2:16: not justified by works of the law but through faith in Christ, Galatians 3:11: no one is justified before God by the law; the righteous shall live by faith, Hebrews 10:38–39: the righteous one shall live by faith, Hebrews 11:6: without faith it is impossible to please God, John 3:16–18: whoever believes has eternal life; whoever does not believe is condemned already, John 5:24: whoever hears and believes has eternal life and does not come into judgement, John 6:28–29: the work of God is that you believe in him whom he has sent, Acts 10:43: everyone who believes in him receives forgiveness of sins, Acts 15:10–11: no yoke on the disciples; we are saved through the grace of the Lord Jesus, Luke 7:50: your faith has saved you, Luke 18:9–14: the tax collector, not the Pharisee who listed his works, went home justified.
\par\medskip\noindent\textbf{\texttt{without\allowbreak{}Hebrews}} --- The canonical witness, without Hebrews’ reading\par\nobreak\noindent
Without Hebrews’ reading.

\begin{derivation}
  \item[(1)] \(\pvar{TestimonyArguments.CanonicalWitness}{1}\)\gloss{Romans 3:28: a person is justified by faith apart from works of the law}
  \item[(2)] \(\pvar{TestimonyArguments.CanonicalWitness}{5}\)\gloss{Ephesians 2:8–9: saved by grace through faith, not of works}
  \item[(3)] \(\pvar{TestimonyArguments.CanonicalWitness}{6}\)\gloss{Titus 3:5: saved not by works done in righteousness but by his mercy}
  \breakrow{(4)} \(\pvar{TestimonyArguments.CanonicalWitness}{2}\)\gloss{Romans 4:4–5: to the one who does not work but believes, faith is counted as righteousness}
  \breakrow{(5)} \(\pvar{TestimonyArguments.CanonicalWitness}{3}\)\gloss{Galatians 2:16: not justified by works of the law but through faith in Christ}
  \breakrow{(6)} \(\pvar{TestimonyArguments.CanonicalWitness}{4}\)\gloss{Galatians 3:11: no one is justified before God by the law; the righteous shall live by faith}
  \breakrow{(7)} \(\pvar{TestimonyArguments.CanonicalWitness}{7}\)\gloss{Hebrews 10:38–39: the righteous one shall live by faith}
  \breakrow{(8)} \(\pvar{TestimonyArguments.CanonicalWitness}{8}\)\gloss{Hebrews 11:6: without faith it is impossible to please God}
  \breakrow{(9)} \(\pvar{TestimonyArguments.CanonicalWitness}{9}\)\gloss{John 3:16–18: whoever believes has eternal life; whoever does not believe is condemned already}
  \breakrow{(10)} \(\pvar{TestimonyArguments.CanonicalWitness}{10}\)\gloss{John 5:24: whoever hears and believes has eternal life and does not come into judgement}
  \breakrow{(11)} \(\pvar{TestimonyArguments.CanonicalWitness}{11}\)\gloss{John 6:28–29: the work of God is that you believe in him whom he has sent}
  \breakrow{(12)} \(\pvar{TestimonyArguments.CanonicalWitness}{12}\)\gloss{Acts 10:43: everyone who believes in him receives forgiveness of sins}
  \breakrow{(13)} \(\pvar{TestimonyArguments.CanonicalWitness}{13}\)\gloss{Acts 15:10–11: no yoke on the disciples; we are saved through the grace of the Lord Jesus}
  \breakrow{(14)} \(\pvar{TestimonyArguments.CanonicalWitness}{14}\)\gloss{Luke 7:50: your faith has saved you}
  \breakrow{(15)} \(\pvar{TestimonyArguments.CanonicalWitness}{15}\)\gloss{Luke 18:9–14: the tax collector, not the Pharisee who listed his works, went home justified}
  \breakrow{(16)} \((\pvar{TestimonyArguments.CanonicalWitness}{1} \land \pvar{TestimonyArguments.CanonicalWitness}{5} \land \pvar{TestimonyArguments.CanonicalWitness}{6}) \rightarrow \pvar{TestimonyArguments.CanonicalWitness}{23}\)\gloss{\emph{would give:} Works are not the ground of justification}
  \breakrow{(17)} \((\pvar{TestimonyArguments.CanonicalWitness}{2} \land \pvar{TestimonyArguments.CanonicalWitness}{3} \land \pvar{TestimonyArguments.CanonicalWitness}{4}) \rightarrow \pvar{TestimonyArguments.CanonicalWitness}{21} \land \pvar{TestimonyArguments.CanonicalWitness}{22}\)\gloss{\emph{would give:} \(\pvar{TestimonyArguments.CanonicalWitness}{21}\) Faith is necessary: no one is saved without it; \(\pvar{TestimonyArguments.CanonicalWitness}{22}\) Faith is sufficient: whoever believes is saved}
  \breakrow{(18)} \((\pvar{TestimonyArguments.CanonicalWitness}{9} \land \pvar{TestimonyArguments.CanonicalWitness}{10} \land \pvar{TestimonyArguments.CanonicalWitness}{11}) \rightarrow \pvar{TestimonyArguments.CanonicalWitness}{22} \land \pvar{TestimonyArguments.CanonicalWitness}{21}\)\gloss{\emph{would give:} \(\pvar{TestimonyArguments.CanonicalWitness}{22}\) Faith is sufficient: whoever believes is saved; \(\pvar{TestimonyArguments.CanonicalWitness}{21}\) Faith is necessary: no one is saved without it}
  \breakrow{(19)} \((\pvar{TestimonyArguments.CanonicalWitness}{12} \land \pvar{TestimonyArguments.CanonicalWitness}{13}) \rightarrow \pvar{TestimonyArguments.CanonicalWitness}{22} \land \pvar{TestimonyArguments.CanonicalWitness}{23}\)\gloss{\emph{would give:} \(\pvar{TestimonyArguments.CanonicalWitness}{22}\) Faith is sufficient: whoever believes is saved; \(\pvar{TestimonyArguments.CanonicalWitness}{23}\) Works are not the ground of justification}
  \breakrow{(20)} \((\pvar{TestimonyArguments.CanonicalWitness}{14} \land \pvar{TestimonyArguments.CanonicalWitness}{15}) \rightarrow \pvar{TestimonyArguments.CanonicalWitness}{22} \land \pvar{TestimonyArguments.CanonicalWitness}{23}\)\gloss{\emph{would give:} \(\pvar{TestimonyArguments.CanonicalWitness}{22}\) Faith is sufficient: whoever believes is saved; \(\pvar{TestimonyArguments.CanonicalWitness}{23}\) Works are not the ground of justification}
  \item[(21)] \((\pvar{TestimonyArguments.CanonicalWitness}{21} \land \pvar{TestimonyArguments.CanonicalWitness}{22} \land \pvar{TestimonyArguments.CanonicalWitness}{23}) \rightarrow \pvar{TestimonyArguments.CanonicalWitness}{24}\)
  \item[\(\vdash\)] \(\pvar{TestimonyArguments.CanonicalWitness}{24}\)\gloss{\emph{at issue:} Justification is by faith alone}
  \keepwithnext
\end{derivation}
\emph{Grounded in scripture alone:} Romans 3:28: a person is justified by faith apart from works of the law, Ephesians 2:8–9: saved by grace through faith, not of works, Titus 3:5: saved not by works done in righteousness but by his mercy, Romans 4:4–5: to the one who does not work but believes, faith is counted as righteousness, Galatians 2:16: not justified by works of the law but through faith in Christ, Galatians 3:11: no one is justified before God by the law; the righteous shall live by faith, Hebrews 10:38–39: the righteous one shall live by faith, Hebrews 11:6: without faith it is impossible to please God, John 3:16–18: whoever believes has eternal life; whoever does not believe is condemned already, John 5:24: whoever hears and believes has eternal life and does not come into judgement, John 6:28–29: the work of God is that you believe in him whom he has sent, Acts 10:43: everyone who believes in him receives forgiveness of sins, Acts 15:10–11: no yoke on the disciples; we are saved through the grace of the Lord Jesus, Luke 7:50: your faith has saved you, Luke 18:9–14: the tax collector, not the Pharisee who listed his works, went home justified.
\par\medskip\noindent\textbf{\texttt{without\allowbreak{}John}} --- The canonical witness, without John’s reading\par\nobreak\noindent
Without John’s reading.

\begin{derivation}
  \item[(1)] \(\pvar{TestimonyArguments.CanonicalWitness}{1}\)\gloss{Romans 3:28: a person is justified by faith apart from works of the law}
  \item[(2)] \(\pvar{TestimonyArguments.CanonicalWitness}{5}\)\gloss{Ephesians 2:8–9: saved by grace through faith, not of works}
  \item[(3)] \(\pvar{TestimonyArguments.CanonicalWitness}{6}\)\gloss{Titus 3:5: saved not by works done in righteousness but by his mercy}
  \breakrow{(4)} \(\pvar{TestimonyArguments.CanonicalWitness}{2}\)\gloss{Romans 4:4–5: to the one who does not work but believes, faith is counted as righteousness}
  \breakrow{(5)} \(\pvar{TestimonyArguments.CanonicalWitness}{3}\)\gloss{Galatians 2:16: not justified by works of the law but through faith in Christ}
  \breakrow{(6)} \(\pvar{TestimonyArguments.CanonicalWitness}{4}\)\gloss{Galatians 3:11: no one is justified before God by the law; the righteous shall live by faith}
  \breakrow{(7)} \(\pvar{TestimonyArguments.CanonicalWitness}{7}\)\gloss{Hebrews 10:38–39: the righteous one shall live by faith}
  \breakrow{(8)} \(\pvar{TestimonyArguments.CanonicalWitness}{8}\)\gloss{Hebrews 11:6: without faith it is impossible to please God}
  \breakrow{(9)} \(\pvar{TestimonyArguments.CanonicalWitness}{9}\)\gloss{John 3:16–18: whoever believes has eternal life; whoever does not believe is condemned already}
  \breakrow{(10)} \(\pvar{TestimonyArguments.CanonicalWitness}{10}\)\gloss{John 5:24: whoever hears and believes has eternal life and does not come into judgement}
  \breakrow{(11)} \(\pvar{TestimonyArguments.CanonicalWitness}{11}\)\gloss{John 6:28–29: the work of God is that you believe in him whom he has sent}
  \breakrow{(12)} \(\pvar{TestimonyArguments.CanonicalWitness}{12}\)\gloss{Acts 10:43: everyone who believes in him receives forgiveness of sins}
  \breakrow{(13)} \(\pvar{TestimonyArguments.CanonicalWitness}{13}\)\gloss{Acts 15:10–11: no yoke on the disciples; we are saved through the grace of the Lord Jesus}
  \breakrow{(14)} \(\pvar{TestimonyArguments.CanonicalWitness}{14}\)\gloss{Luke 7:50: your faith has saved you}
  \breakrow{(15)} \(\pvar{TestimonyArguments.CanonicalWitness}{15}\)\gloss{Luke 18:9–14: the tax collector, not the Pharisee who listed his works, went home justified}
  \breakrow{(16)} \((\pvar{TestimonyArguments.CanonicalWitness}{1} \land \pvar{TestimonyArguments.CanonicalWitness}{5} \land \pvar{TestimonyArguments.CanonicalWitness}{6}) \rightarrow \pvar{TestimonyArguments.CanonicalWitness}{23}\)\gloss{\emph{would give:} Works are not the ground of justification}
  \breakrow{(17)} \((\pvar{TestimonyArguments.CanonicalWitness}{2} \land \pvar{TestimonyArguments.CanonicalWitness}{3} \land \pvar{TestimonyArguments.CanonicalWitness}{4}) \rightarrow \pvar{TestimonyArguments.CanonicalWitness}{21} \land \pvar{TestimonyArguments.CanonicalWitness}{22}\)\gloss{\emph{would give:} \(\pvar{TestimonyArguments.CanonicalWitness}{21}\) Faith is necessary: no one is saved without it; \(\pvar{TestimonyArguments.CanonicalWitness}{22}\) Faith is sufficient: whoever believes is saved}
  \breakrow{(18)} \((\pvar{TestimonyArguments.CanonicalWitness}{7} \land \pvar{TestimonyArguments.CanonicalWitness}{8}) \rightarrow \pvar{TestimonyArguments.CanonicalWitness}{21}\)\gloss{\emph{would give:} Faith is necessary: no one is saved without it}
  \breakrow{(19)} \((\pvar{TestimonyArguments.CanonicalWitness}{12} \land \pvar{TestimonyArguments.CanonicalWitness}{13}) \rightarrow \pvar{TestimonyArguments.CanonicalWitness}{22} \land \pvar{TestimonyArguments.CanonicalWitness}{23}\)\gloss{\emph{would give:} \(\pvar{TestimonyArguments.CanonicalWitness}{22}\) Faith is sufficient: whoever believes is saved; \(\pvar{TestimonyArguments.CanonicalWitness}{23}\) Works are not the ground of justification}
  \breakrow{(20)} \((\pvar{TestimonyArguments.CanonicalWitness}{14} \land \pvar{TestimonyArguments.CanonicalWitness}{15}) \rightarrow \pvar{TestimonyArguments.CanonicalWitness}{22} \land \pvar{TestimonyArguments.CanonicalWitness}{23}\)\gloss{\emph{would give:} \(\pvar{TestimonyArguments.CanonicalWitness}{22}\) Faith is sufficient: whoever believes is saved; \(\pvar{TestimonyArguments.CanonicalWitness}{23}\) Works are not the ground of justification}
  \item[(21)] \((\pvar{TestimonyArguments.CanonicalWitness}{21} \land \pvar{TestimonyArguments.CanonicalWitness}{22} \land \pvar{TestimonyArguments.CanonicalWitness}{23}) \rightarrow \pvar{TestimonyArguments.CanonicalWitness}{24}\)
  \item[\(\vdash\)] \(\pvar{TestimonyArguments.CanonicalWitness}{24}\)\gloss{\emph{at issue:} Justification is by faith alone}
  \keepwithnext
\end{derivation}
\emph{Grounded in scripture alone:} Romans 3:28: a person is justified by faith apart from works of the law, Ephesians 2:8–9: saved by grace through faith, not of works, Titus 3:5: saved not by works done in righteousness but by his mercy, Romans 4:4–5: to the one who does not work but believes, faith is counted as righteousness, Galatians 2:16: not justified by works of the law but through faith in Christ, Galatians 3:11: no one is justified before God by the law; the righteous shall live by faith, Hebrews 10:38–39: the righteous one shall live by faith, Hebrews 11:6: without faith it is impossible to please God, John 3:16–18: whoever believes has eternal life; whoever does not believe is condemned already, John 5:24: whoever hears and believes has eternal life and does not come into judgement, John 6:28–29: the work of God is that you believe in him whom he has sent, Acts 10:43: everyone who believes in him receives forgiveness of sins, Acts 15:10–11: no yoke on the disciples; we are saved through the grace of the Lord Jesus, Luke 7:50: your faith has saved you, Luke 18:9–14: the tax collector, not the Pharisee who listed his works, went home justified.
\par\medskip\noindent\textbf{\texttt{without\allowbreak{}Peter}} --- The canonical witness, without Peter’s reading\par\nobreak\noindent
Without Peter’s reading.

\begin{derivation}
  \item[(1)] \(\pvar{TestimonyArguments.CanonicalWitness}{1}\)\gloss{Romans 3:28: a person is justified by faith apart from works of the law}
  \item[(2)] \(\pvar{TestimonyArguments.CanonicalWitness}{5}\)\gloss{Ephesians 2:8–9: saved by grace through faith, not of works}
  \item[(3)] \(\pvar{TestimonyArguments.CanonicalWitness}{6}\)\gloss{Titus 3:5: saved not by works done in righteousness but by his mercy}
  \breakrow{(4)} \(\pvar{TestimonyArguments.CanonicalWitness}{2}\)\gloss{Romans 4:4–5: to the one who does not work but believes, faith is counted as righteousness}
  \breakrow{(5)} \(\pvar{TestimonyArguments.CanonicalWitness}{3}\)\gloss{Galatians 2:16: not justified by works of the law but through faith in Christ}
  \breakrow{(6)} \(\pvar{TestimonyArguments.CanonicalWitness}{4}\)\gloss{Galatians 3:11: no one is justified before God by the law; the righteous shall live by faith}
  \breakrow{(7)} \(\pvar{TestimonyArguments.CanonicalWitness}{7}\)\gloss{Hebrews 10:38–39: the righteous one shall live by faith}
  \breakrow{(8)} \(\pvar{TestimonyArguments.CanonicalWitness}{8}\)\gloss{Hebrews 11:6: without faith it is impossible to please God}
  \breakrow{(9)} \(\pvar{TestimonyArguments.CanonicalWitness}{9}\)\gloss{John 3:16–18: whoever believes has eternal life; whoever does not believe is condemned already}
  \breakrow{(10)} \(\pvar{TestimonyArguments.CanonicalWitness}{10}\)\gloss{John 5:24: whoever hears and believes has eternal life and does not come into judgement}
  \breakrow{(11)} \(\pvar{TestimonyArguments.CanonicalWitness}{11}\)\gloss{John 6:28–29: the work of God is that you believe in him whom he has sent}
  \breakrow{(12)} \(\pvar{TestimonyArguments.CanonicalWitness}{12}\)\gloss{Acts 10:43: everyone who believes in him receives forgiveness of sins}
  \breakrow{(13)} \(\pvar{TestimonyArguments.CanonicalWitness}{13}\)\gloss{Acts 15:10–11: no yoke on the disciples; we are saved through the grace of the Lord Jesus}
  \breakrow{(14)} \(\pvar{TestimonyArguments.CanonicalWitness}{14}\)\gloss{Luke 7:50: your faith has saved you}
  \breakrow{(15)} \(\pvar{TestimonyArguments.CanonicalWitness}{15}\)\gloss{Luke 18:9–14: the tax collector, not the Pharisee who listed his works, went home justified}
  \breakrow{(16)} \((\pvar{TestimonyArguments.CanonicalWitness}{1} \land \pvar{TestimonyArguments.CanonicalWitness}{5} \land \pvar{TestimonyArguments.CanonicalWitness}{6}) \rightarrow \pvar{TestimonyArguments.CanonicalWitness}{23}\)\gloss{\emph{would give:} Works are not the ground of justification}
  \breakrow{(17)} \((\pvar{TestimonyArguments.CanonicalWitness}{2} \land \pvar{TestimonyArguments.CanonicalWitness}{3} \land \pvar{TestimonyArguments.CanonicalWitness}{4}) \rightarrow \pvar{TestimonyArguments.CanonicalWitness}{21} \land \pvar{TestimonyArguments.CanonicalWitness}{22}\)\gloss{\emph{would give:} \(\pvar{TestimonyArguments.CanonicalWitness}{21}\) Faith is necessary: no one is saved without it; \(\pvar{TestimonyArguments.CanonicalWitness}{22}\) Faith is sufficient: whoever believes is saved}
  \breakrow{(18)} \((\pvar{TestimonyArguments.CanonicalWitness}{7} \land \pvar{TestimonyArguments.CanonicalWitness}{8}) \rightarrow \pvar{TestimonyArguments.CanonicalWitness}{21}\)\gloss{\emph{would give:} Faith is necessary: no one is saved without it}
  \breakrow{(19)} \((\pvar{TestimonyArguments.CanonicalWitness}{9} \land \pvar{TestimonyArguments.CanonicalWitness}{10} \land \pvar{TestimonyArguments.CanonicalWitness}{11}) \rightarrow \pvar{TestimonyArguments.CanonicalWitness}{22} \land \pvar{TestimonyArguments.CanonicalWitness}{21}\)\gloss{\emph{would give:} \(\pvar{TestimonyArguments.CanonicalWitness}{22}\) Faith is sufficient: whoever believes is saved; \(\pvar{TestimonyArguments.CanonicalWitness}{21}\) Faith is necessary: no one is saved without it}
  \breakrow{(20)} \((\pvar{TestimonyArguments.CanonicalWitness}{14} \land \pvar{TestimonyArguments.CanonicalWitness}{15}) \rightarrow \pvar{TestimonyArguments.CanonicalWitness}{22} \land \pvar{TestimonyArguments.CanonicalWitness}{23}\)\gloss{\emph{would give:} \(\pvar{TestimonyArguments.CanonicalWitness}{22}\) Faith is sufficient: whoever believes is saved; \(\pvar{TestimonyArguments.CanonicalWitness}{23}\) Works are not the ground of justification}
  \item[(21)] \((\pvar{TestimonyArguments.CanonicalWitness}{21} \land \pvar{TestimonyArguments.CanonicalWitness}{22} \land \pvar{TestimonyArguments.CanonicalWitness}{23}) \rightarrow \pvar{TestimonyArguments.CanonicalWitness}{24}\)
  \item[\(\vdash\)] \(\pvar{TestimonyArguments.CanonicalWitness}{24}\)\gloss{\emph{at issue:} Justification is by faith alone}
  \keepwithnext
\end{derivation}
\emph{Grounded in scripture alone:} Romans 3:28: a person is justified by faith apart from works of the law, Ephesians 2:8–9: saved by grace through faith, not of works, Titus 3:5: saved not by works done in righteousness but by his mercy, Romans 4:4–5: to the one who does not work but believes, faith is counted as righteousness, Galatians 2:16: not justified by works of the law but through faith in Christ, Galatians 3:11: no one is justified before God by the law; the righteous shall live by faith, Hebrews 10:38–39: the righteous one shall live by faith, Hebrews 11:6: without faith it is impossible to please God, John 3:16–18: whoever believes has eternal life; whoever does not believe is condemned already, John 5:24: whoever hears and believes has eternal life and does not come into judgement, John 6:28–29: the work of God is that you believe in him whom he has sent, Acts 10:43: everyone who believes in him receives forgiveness of sins, Acts 15:10–11: no yoke on the disciples; we are saved through the grace of the Lord Jesus, Luke 7:50: your faith has saved you, Luke 18:9–14: the tax collector, not the Pharisee who listed his works, went home justified.
\par\medskip\noindent\textbf{\texttt{without\allowbreak{}Jesus}} --- The canonical witness, without the reading of Jesus’ words\par\nobreak\noindent
Without the reading of Jesus’ words.

\begin{derivation}
  \item[(1)] \(\pvar{TestimonyArguments.CanonicalWitness}{1}\)\gloss{Romans 3:28: a person is justified by faith apart from works of the law}
  \item[(2)] \(\pvar{TestimonyArguments.CanonicalWitness}{5}\)\gloss{Ephesians 2:8–9: saved by grace through faith, not of works}
  \item[(3)] \(\pvar{TestimonyArguments.CanonicalWitness}{6}\)\gloss{Titus 3:5: saved not by works done in righteousness but by his mercy}
  \breakrow{(4)} \(\pvar{TestimonyArguments.CanonicalWitness}{2}\)\gloss{Romans 4:4–5: to the one who does not work but believes, faith is counted as righteousness}
  \breakrow{(5)} \(\pvar{TestimonyArguments.CanonicalWitness}{3}\)\gloss{Galatians 2:16: not justified by works of the law but through faith in Christ}
  \breakrow{(6)} \(\pvar{TestimonyArguments.CanonicalWitness}{4}\)\gloss{Galatians 3:11: no one is justified before God by the law; the righteous shall live by faith}
  \breakrow{(7)} \(\pvar{TestimonyArguments.CanonicalWitness}{7}\)\gloss{Hebrews 10:38–39: the righteous one shall live by faith}
  \breakrow{(8)} \(\pvar{TestimonyArguments.CanonicalWitness}{8}\)\gloss{Hebrews 11:6: without faith it is impossible to please God}
  \breakrow{(9)} \(\pvar{TestimonyArguments.CanonicalWitness}{9}\)\gloss{John 3:16–18: whoever believes has eternal life; whoever does not believe is condemned already}
  \breakrow{(10)} \(\pvar{TestimonyArguments.CanonicalWitness}{10}\)\gloss{John 5:24: whoever hears and believes has eternal life and does not come into judgement}
  \breakrow{(11)} \(\pvar{TestimonyArguments.CanonicalWitness}{11}\)\gloss{John 6:28–29: the work of God is that you believe in him whom he has sent}
  \breakrow{(12)} \(\pvar{TestimonyArguments.CanonicalWitness}{12}\)\gloss{Acts 10:43: everyone who believes in him receives forgiveness of sins}
  \breakrow{(13)} \(\pvar{TestimonyArguments.CanonicalWitness}{13}\)\gloss{Acts 15:10–11: no yoke on the disciples; we are saved through the grace of the Lord Jesus}
  \breakrow{(14)} \(\pvar{TestimonyArguments.CanonicalWitness}{14}\)\gloss{Luke 7:50: your faith has saved you}
  \breakrow{(15)} \(\pvar{TestimonyArguments.CanonicalWitness}{15}\)\gloss{Luke 18:9–14: the tax collector, not the Pharisee who listed his works, went home justified}
  \breakrow{(16)} \((\pvar{TestimonyArguments.CanonicalWitness}{1} \land \pvar{TestimonyArguments.CanonicalWitness}{5} \land \pvar{TestimonyArguments.CanonicalWitness}{6}) \rightarrow \pvar{TestimonyArguments.CanonicalWitness}{23}\)\gloss{\emph{would give:} Works are not the ground of justification}
  \breakrow{(17)} \((\pvar{TestimonyArguments.CanonicalWitness}{2} \land \pvar{TestimonyArguments.CanonicalWitness}{3} \land \pvar{TestimonyArguments.CanonicalWitness}{4}) \rightarrow \pvar{TestimonyArguments.CanonicalWitness}{21} \land \pvar{TestimonyArguments.CanonicalWitness}{22}\)\gloss{\emph{would give:} \(\pvar{TestimonyArguments.CanonicalWitness}{21}\) Faith is necessary: no one is saved without it; \(\pvar{TestimonyArguments.CanonicalWitness}{22}\) Faith is sufficient: whoever believes is saved}
  \breakrow{(18)} \((\pvar{TestimonyArguments.CanonicalWitness}{7} \land \pvar{TestimonyArguments.CanonicalWitness}{8}) \rightarrow \pvar{TestimonyArguments.CanonicalWitness}{21}\)\gloss{\emph{would give:} Faith is necessary: no one is saved without it}
  \breakrow{(19)} \((\pvar{TestimonyArguments.CanonicalWitness}{9} \land \pvar{TestimonyArguments.CanonicalWitness}{10} \land \pvar{TestimonyArguments.CanonicalWitness}{11}) \rightarrow \pvar{TestimonyArguments.CanonicalWitness}{22} \land \pvar{TestimonyArguments.CanonicalWitness}{21}\)\gloss{\emph{would give:} \(\pvar{TestimonyArguments.CanonicalWitness}{22}\) Faith is sufficient: whoever believes is saved; \(\pvar{TestimonyArguments.CanonicalWitness}{21}\) Faith is necessary: no one is saved without it}
  \breakrow{(20)} \((\pvar{TestimonyArguments.CanonicalWitness}{12} \land \pvar{TestimonyArguments.CanonicalWitness}{13}) \rightarrow \pvar{TestimonyArguments.CanonicalWitness}{22} \land \pvar{TestimonyArguments.CanonicalWitness}{23}\)\gloss{\emph{would give:} \(\pvar{TestimonyArguments.CanonicalWitness}{22}\) Faith is sufficient: whoever believes is saved; \(\pvar{TestimonyArguments.CanonicalWitness}{23}\) Works are not the ground of justification}
  \item[(21)] \((\pvar{TestimonyArguments.CanonicalWitness}{21} \land \pvar{TestimonyArguments.CanonicalWitness}{22} \land \pvar{TestimonyArguments.CanonicalWitness}{23}) \rightarrow \pvar{TestimonyArguments.CanonicalWitness}{24}\)
  \item[\(\vdash\)] \(\pvar{TestimonyArguments.CanonicalWitness}{24}\)\gloss{\emph{at issue:} Justification is by faith alone}
  \keepwithnext
\end{derivation}
\emph{Grounded in scripture alone:} Romans 3:28: a person is justified by faith apart from works of the law, Ephesians 2:8–9: saved by grace through faith, not of works, Titus 3:5: saved not by works done in righteousness but by his mercy, Romans 4:4–5: to the one who does not work but believes, faith is counted as righteousness, Galatians 2:16: not justified by works of the law but through faith in Christ, Galatians 3:11: no one is justified before God by the law; the righteous shall live by faith, Hebrews 10:38–39: the righteous one shall live by faith, Hebrews 11:6: without faith it is impossible to please God, John 3:16–18: whoever believes has eternal life; whoever does not believe is condemned already, John 5:24: whoever hears and believes has eternal life and does not come into judgement, John 6:28–29: the work of God is that you believe in him whom he has sent, Acts 10:43: everyone who believes in him receives forgiveness of sins, Acts 15:10–11: no yoke on the disciples; we are saved through the grace of the Lord Jesus, Luke 7:50: your faith has saved you, Luke 18:9–14: the tax collector, not the Pharisee who listed his works, went home justified.
\paragraph*{The opponent’s two ways}
\leavevmode\par\nobreak
\par\medskip\noindent\textbf{\texttt{without\allowbreak{}The\allowbreak{}Works\allowbreak{}Witnesses}} --- The canonical witness, without Paul’s on works, Peter’s and Jesus’ readings\par\nobreak\noindent
Without the three readings that say works are not the ground — Paul’s on works, Peter’s and Jesus’.

\begin{derivation}
  \item[(1)] \(\pvar{TestimonyArguments.CanonicalWitness}{1}\)\gloss{Romans 3:28: a person is justified by faith apart from works of the law}
  \item[(2)] \(\pvar{TestimonyArguments.CanonicalWitness}{5}\)\gloss{Ephesians 2:8–9: saved by grace through faith, not of works}
  \item[(3)] \(\pvar{TestimonyArguments.CanonicalWitness}{6}\)\gloss{Titus 3:5: saved not by works done in righteousness but by his mercy}
  \breakrow{(4)} \(\pvar{TestimonyArguments.CanonicalWitness}{2}\)\gloss{Romans 4:4–5: to the one who does not work but believes, faith is counted as righteousness}
  \breakrow{(5)} \(\pvar{TestimonyArguments.CanonicalWitness}{3}\)\gloss{Galatians 2:16: not justified by works of the law but through faith in Christ}
  \breakrow{(6)} \(\pvar{TestimonyArguments.CanonicalWitness}{4}\)\gloss{Galatians 3:11: no one is justified before God by the law; the righteous shall live by faith}
  \breakrow{(7)} \(\pvar{TestimonyArguments.CanonicalWitness}{7}\)\gloss{Hebrews 10:38–39: the righteous one shall live by faith}
  \breakrow{(8)} \(\pvar{TestimonyArguments.CanonicalWitness}{8}\)\gloss{Hebrews 11:6: without faith it is impossible to please God}
  \breakrow{(9)} \(\pvar{TestimonyArguments.CanonicalWitness}{9}\)\gloss{John 3:16–18: whoever believes has eternal life; whoever does not believe is condemned already}
  \breakrow{(10)} \(\pvar{TestimonyArguments.CanonicalWitness}{10}\)\gloss{John 5:24: whoever hears and believes has eternal life and does not come into judgement}
  \breakrow{(11)} \(\pvar{TestimonyArguments.CanonicalWitness}{11}\)\gloss{John 6:28–29: the work of God is that you believe in him whom he has sent}
  \breakrow{(12)} \(\pvar{TestimonyArguments.CanonicalWitness}{12}\)\gloss{Acts 10:43: everyone who believes in him receives forgiveness of sins}
  \breakrow{(13)} \(\pvar{TestimonyArguments.CanonicalWitness}{13}\)\gloss{Acts 15:10–11: no yoke on the disciples; we are saved through the grace of the Lord Jesus}
  \breakrow{(14)} \(\pvar{TestimonyArguments.CanonicalWitness}{14}\)\gloss{Luke 7:50: your faith has saved you}
  \breakrow{(15)} \(\pvar{TestimonyArguments.CanonicalWitness}{15}\)\gloss{Luke 18:9–14: the tax collector, not the Pharisee who listed his works, went home justified}
  \breakrow{(16)} \((\pvar{TestimonyArguments.CanonicalWitness}{2} \land \pvar{TestimonyArguments.CanonicalWitness}{3} \land \pvar{TestimonyArguments.CanonicalWitness}{4}) \rightarrow \pvar{TestimonyArguments.CanonicalWitness}{21} \land \pvar{TestimonyArguments.CanonicalWitness}{22}\)\gloss{\emph{would give:} \(\pvar{TestimonyArguments.CanonicalWitness}{21}\) Faith is necessary: no one is saved without it; \(\pvar{TestimonyArguments.CanonicalWitness}{22}\) Faith is sufficient: whoever believes is saved}
  \breakrow{(17)} \((\pvar{TestimonyArguments.CanonicalWitness}{7} \land \pvar{TestimonyArguments.CanonicalWitness}{8}) \rightarrow \pvar{TestimonyArguments.CanonicalWitness}{21}\)\gloss{\emph{would give:} Faith is necessary: no one is saved without it}
  \breakrow{(18)} \((\pvar{TestimonyArguments.CanonicalWitness}{9} \land \pvar{TestimonyArguments.CanonicalWitness}{10} \land \pvar{TestimonyArguments.CanonicalWitness}{11}) \rightarrow \pvar{TestimonyArguments.CanonicalWitness}{22} \land \pvar{TestimonyArguments.CanonicalWitness}{21}\)\gloss{\emph{would give:} \(\pvar{TestimonyArguments.CanonicalWitness}{22}\) Faith is sufficient: whoever believes is saved; \(\pvar{TestimonyArguments.CanonicalWitness}{21}\) Faith is necessary: no one is saved without it}
  \item[(19)] \((\pvar{TestimonyArguments.CanonicalWitness}{21} \land \pvar{TestimonyArguments.CanonicalWitness}{22} \land \pvar{TestimonyArguments.CanonicalWitness}{23}) \rightarrow \pvar{TestimonyArguments.CanonicalWitness}{24}\)
  \item[\(\vdash\)] \(\pvar{TestimonyArguments.CanonicalWitness}{24}\)\gloss{\emph{at issue:} Justification is by faith alone}
  \keepwithnext
\end{derivation}
\emph{Grounded in scripture alone:} Romans 3:28: a person is justified by faith apart from works of the law, Ephesians 2:8–9: saved by grace through faith, not of works, Titus 3:5: saved not by works done in righteousness but by his mercy, Romans 4:4–5: to the one who does not work but believes, faith is counted as righteousness, Galatians 2:16: not justified by works of the law but through faith in Christ, Galatians 3:11: no one is justified before God by the law; the righteous shall live by faith, Hebrews 10:38–39: the righteous one shall live by faith, Hebrews 11:6: without faith it is impossible to please God, John 3:16–18: whoever believes has eternal life; whoever does not believe is condemned already, John 5:24: whoever hears and believes has eternal life and does not come into judgement, John 6:28–29: the work of God is that you believe in him whom he has sent, Acts 10:43: everyone who believes in him receives forgiveness of sins, Acts 15:10–11: no yoke on the disciples; we are saved through the grace of the Lord Jesus, Luke 7:50: your faith has saved you, Luke 18:9–14: the tax collector, not the Pharisee who listed his works, went home justified.
\par\medskip\noindent\textbf{\texttt{without\allowbreak{}The\allowbreak{}Necessity\allowbreak{}Witnesses}} --- The canonical witness, without Paul’s on faith, Hebrews’ and John’s readings\par\nobreak\noindent
Without the three readings that say faith is necessary — Paul’s on faith, Hebrews’ and John’s.

\begin{derivation}
  \item[(1)] \(\pvar{TestimonyArguments.CanonicalWitness}{1}\)\gloss{Romans 3:28: a person is justified by faith apart from works of the law}
  \item[(2)] \(\pvar{TestimonyArguments.CanonicalWitness}{5}\)\gloss{Ephesians 2:8–9: saved by grace through faith, not of works}
  \item[(3)] \(\pvar{TestimonyArguments.CanonicalWitness}{6}\)\gloss{Titus 3:5: saved not by works done in righteousness but by his mercy}
  \breakrow{(4)} \(\pvar{TestimonyArguments.CanonicalWitness}{2}\)\gloss{Romans 4:4–5: to the one who does not work but believes, faith is counted as righteousness}
  \breakrow{(5)} \(\pvar{TestimonyArguments.CanonicalWitness}{3}\)\gloss{Galatians 2:16: not justified by works of the law but through faith in Christ}
  \breakrow{(6)} \(\pvar{TestimonyArguments.CanonicalWitness}{4}\)\gloss{Galatians 3:11: no one is justified before God by the law; the righteous shall live by faith}
  \breakrow{(7)} \(\pvar{TestimonyArguments.CanonicalWitness}{7}\)\gloss{Hebrews 10:38–39: the righteous one shall live by faith}
  \breakrow{(8)} \(\pvar{TestimonyArguments.CanonicalWitness}{8}\)\gloss{Hebrews 11:6: without faith it is impossible to please God}
  \breakrow{(9)} \(\pvar{TestimonyArguments.CanonicalWitness}{9}\)\gloss{John 3:16–18: whoever believes has eternal life; whoever does not believe is condemned already}
  \breakrow{(10)} \(\pvar{TestimonyArguments.CanonicalWitness}{10}\)\gloss{John 5:24: whoever hears and believes has eternal life and does not come into judgement}
  \breakrow{(11)} \(\pvar{TestimonyArguments.CanonicalWitness}{11}\)\gloss{John 6:28–29: the work of God is that you believe in him whom he has sent}
  \breakrow{(12)} \(\pvar{TestimonyArguments.CanonicalWitness}{12}\)\gloss{Acts 10:43: everyone who believes in him receives forgiveness of sins}
  \breakrow{(13)} \(\pvar{TestimonyArguments.CanonicalWitness}{13}\)\gloss{Acts 15:10–11: no yoke on the disciples; we are saved through the grace of the Lord Jesus}
  \breakrow{(14)} \(\pvar{TestimonyArguments.CanonicalWitness}{14}\)\gloss{Luke 7:50: your faith has saved you}
  \breakrow{(15)} \(\pvar{TestimonyArguments.CanonicalWitness}{15}\)\gloss{Luke 18:9–14: the tax collector, not the Pharisee who listed his works, went home justified}
  \breakrow{(16)} \((\pvar{TestimonyArguments.CanonicalWitness}{1} \land \pvar{TestimonyArguments.CanonicalWitness}{5} \land \pvar{TestimonyArguments.CanonicalWitness}{6}) \rightarrow \pvar{TestimonyArguments.CanonicalWitness}{23}\)\gloss{\emph{would give:} Works are not the ground of justification}
  \breakrow{(17)} \((\pvar{TestimonyArguments.CanonicalWitness}{12} \land \pvar{TestimonyArguments.CanonicalWitness}{13}) \rightarrow \pvar{TestimonyArguments.CanonicalWitness}{22} \land \pvar{TestimonyArguments.CanonicalWitness}{23}\)\gloss{\emph{would give:} \(\pvar{TestimonyArguments.CanonicalWitness}{22}\) Faith is sufficient: whoever believes is saved; \(\pvar{TestimonyArguments.CanonicalWitness}{23}\) Works are not the ground of justification}
  \breakrow{(18)} \((\pvar{TestimonyArguments.CanonicalWitness}{14} \land \pvar{TestimonyArguments.CanonicalWitness}{15}) \rightarrow \pvar{TestimonyArguments.CanonicalWitness}{22} \land \pvar{TestimonyArguments.CanonicalWitness}{23}\)\gloss{\emph{would give:} \(\pvar{TestimonyArguments.CanonicalWitness}{22}\) Faith is sufficient: whoever believes is saved; \(\pvar{TestimonyArguments.CanonicalWitness}{23}\) Works are not the ground of justification}
  \item[(19)] \((\pvar{TestimonyArguments.CanonicalWitness}{21} \land \pvar{TestimonyArguments.CanonicalWitness}{22} \land \pvar{TestimonyArguments.CanonicalWitness}{23}) \rightarrow \pvar{TestimonyArguments.CanonicalWitness}{24}\)
  \item[\(\vdash\)] \(\pvar{TestimonyArguments.CanonicalWitness}{24}\)\gloss{\emph{at issue:} Justification is by faith alone}
  \keepwithnext
\end{derivation}
\emph{Grounded in scripture alone:} Romans 3:28: a person is justified by faith apart from works of the law, Ephesians 2:8–9: saved by grace through faith, not of works, Titus 3:5: saved not by works done in righteousness but by his mercy, Romans 4:4–5: to the one who does not work but believes, faith is counted as righteousness, Galatians 2:16: not justified by works of the law but through faith in Christ, Galatians 3:11: no one is justified before God by the law; the righteous shall live by faith, Hebrews 10:38–39: the righteous one shall live by faith, Hebrews 11:6: without faith it is impossible to please God, John 3:16–18: whoever believes has eternal life; whoever does not believe is condemned already, John 5:24: whoever hears and believes has eternal life and does not come into judgement, John 6:28–29: the work of God is that you believe in him whom he has sent, Acts 10:43: everyone who believes in him receives forgiveness of sins, Acts 15:10–11: no yoke on the disciples; we are saved through the grace of the Lord Jesus, Luke 7:50: your faith has saved you, Luke 18:9–14: the tax collector, not the Pharisee who listed his works, went home justified.
\par\medskip\noindent\textbf{\texttt{without\allowbreak{}The\allowbreak{}Sufficiency\allowbreak{}Witnesses}} --- The canonical witness, without Paul’s on faith, John’s, Peter’s and Jesus’ readings\par\nobreak\noindent
Without the four readings that say faith is sufficient — Paul’s on faith, John’s, Peter’s and Jesus’.

\begin{derivation}
  \item[(1)] \(\pvar{TestimonyArguments.CanonicalWitness}{1}\)\gloss{Romans 3:28: a person is justified by faith apart from works of the law}
  \item[(2)] \(\pvar{TestimonyArguments.CanonicalWitness}{5}\)\gloss{Ephesians 2:8–9: saved by grace through faith, not of works}
  \item[(3)] \(\pvar{TestimonyArguments.CanonicalWitness}{6}\)\gloss{Titus 3:5: saved not by works done in righteousness but by his mercy}
  \breakrow{(4)} \(\pvar{TestimonyArguments.CanonicalWitness}{2}\)\gloss{Romans 4:4–5: to the one who does not work but believes, faith is counted as righteousness}
  \breakrow{(5)} \(\pvar{TestimonyArguments.CanonicalWitness}{3}\)\gloss{Galatians 2:16: not justified by works of the law but through faith in Christ}
  \breakrow{(6)} \(\pvar{TestimonyArguments.CanonicalWitness}{4}\)\gloss{Galatians 3:11: no one is justified before God by the law; the righteous shall live by faith}
  \breakrow{(7)} \(\pvar{TestimonyArguments.CanonicalWitness}{7}\)\gloss{Hebrews 10:38–39: the righteous one shall live by faith}
  \breakrow{(8)} \(\pvar{TestimonyArguments.CanonicalWitness}{8}\)\gloss{Hebrews 11:6: without faith it is impossible to please God}
  \breakrow{(9)} \(\pvar{TestimonyArguments.CanonicalWitness}{9}\)\gloss{John 3:16–18: whoever believes has eternal life; whoever does not believe is condemned already}
  \breakrow{(10)} \(\pvar{TestimonyArguments.CanonicalWitness}{10}\)\gloss{John 5:24: whoever hears and believes has eternal life and does not come into judgement}
  \breakrow{(11)} \(\pvar{TestimonyArguments.CanonicalWitness}{11}\)\gloss{John 6:28–29: the work of God is that you believe in him whom he has sent}
  \breakrow{(12)} \(\pvar{TestimonyArguments.CanonicalWitness}{12}\)\gloss{Acts 10:43: everyone who believes in him receives forgiveness of sins}
  \breakrow{(13)} \(\pvar{TestimonyArguments.CanonicalWitness}{13}\)\gloss{Acts 15:10–11: no yoke on the disciples; we are saved through the grace of the Lord Jesus}
  \breakrow{(14)} \(\pvar{TestimonyArguments.CanonicalWitness}{14}\)\gloss{Luke 7:50: your faith has saved you}
  \breakrow{(15)} \(\pvar{TestimonyArguments.CanonicalWitness}{15}\)\gloss{Luke 18:9–14: the tax collector, not the Pharisee who listed his works, went home justified}
  \breakrow{(16)} \((\pvar{TestimonyArguments.CanonicalWitness}{1} \land \pvar{TestimonyArguments.CanonicalWitness}{5} \land \pvar{TestimonyArguments.CanonicalWitness}{6}) \rightarrow \pvar{TestimonyArguments.CanonicalWitness}{23}\)\gloss{\emph{would give:} Works are not the ground of justification}
  \breakrow{(17)} \((\pvar{TestimonyArguments.CanonicalWitness}{7} \land \pvar{TestimonyArguments.CanonicalWitness}{8}) \rightarrow \pvar{TestimonyArguments.CanonicalWitness}{21}\)\gloss{\emph{would give:} Faith is necessary: no one is saved without it}
  \item[(18)] \((\pvar{TestimonyArguments.CanonicalWitness}{21} \land \pvar{TestimonyArguments.CanonicalWitness}{22} \land \pvar{TestimonyArguments.CanonicalWitness}{23}) \rightarrow \pvar{TestimonyArguments.CanonicalWitness}{24}\)
  \item[\(\vdash\)] \(\pvar{TestimonyArguments.CanonicalWitness}{24}\)\gloss{\emph{at issue:} Justification is by faith alone}
  \keepwithnext
\end{derivation}
\emph{Grounded in scripture alone:} Romans 3:28: a person is justified by faith apart from works of the law, Ephesians 2:8–9: saved by grace through faith, not of works, Titus 3:5: saved not by works done in righteousness but by his mercy, Romans 4:4–5: to the one who does not work but believes, faith is counted as righteousness, Galatians 2:16: not justified by works of the law but through faith in Christ, Galatians 3:11: no one is justified before God by the law; the righteous shall live by faith, Hebrews 10:38–39: the righteous one shall live by faith, Hebrews 11:6: without faith it is impossible to please God, John 3:16–18: whoever believes has eternal life; whoever does not believe is condemned already, John 5:24: whoever hears and believes has eternal life and does not come into judgement, John 6:28–29: the work of God is that you believe in him whom he has sent, Acts 10:43: everyone who believes in him receives forgiveness of sins, Acts 15:10–11: no yoke on the disciples; we are saved through the grace of the Lord Jesus, Luke 7:50: your faith has saved you, Luke 18:9–14: the tax collector, not the Pharisee who listed his works, went home justified.
\par\medskip\noindent\textbf{\texttt{hebrews\allowbreak{}And\allowbreak{}Jesus}} --- Hebrews and Jesus’ words alone\par\nobreak\noindent
Two corpora alone: Hebrews and Jesus’ words, with Paul’s, John’s and Peter’s readings rejected.

\begin{derivation}
  \item[(1)] \(\pvar{TestimonyArguments.CanonicalWitness}{1}\)\gloss{Romans 3:28: a person is justified by faith apart from works of the law}
  \item[(2)] \(\pvar{TestimonyArguments.CanonicalWitness}{5}\)\gloss{Ephesians 2:8–9: saved by grace through faith, not of works}
  \item[(3)] \(\pvar{TestimonyArguments.CanonicalWitness}{6}\)\gloss{Titus 3:5: saved not by works done in righteousness but by his mercy}
  \breakrow{(4)} \(\pvar{TestimonyArguments.CanonicalWitness}{2}\)\gloss{Romans 4:4–5: to the one who does not work but believes, faith is counted as righteousness}
  \breakrow{(5)} \(\pvar{TestimonyArguments.CanonicalWitness}{3}\)\gloss{Galatians 2:16: not justified by works of the law but through faith in Christ}
  \breakrow{(6)} \(\pvar{TestimonyArguments.CanonicalWitness}{4}\)\gloss{Galatians 3:11: no one is justified before God by the law; the righteous shall live by faith}
  \breakrow{(7)} \(\pvar{TestimonyArguments.CanonicalWitness}{7}\)\gloss{Hebrews 10:38–39: the righteous one shall live by faith}
  \breakrow{(8)} \(\pvar{TestimonyArguments.CanonicalWitness}{8}\)\gloss{Hebrews 11:6: without faith it is impossible to please God}
  \breakrow{(9)} \(\pvar{TestimonyArguments.CanonicalWitness}{9}\)\gloss{John 3:16–18: whoever believes has eternal life; whoever does not believe is condemned already}
  \breakrow{(10)} \(\pvar{TestimonyArguments.CanonicalWitness}{10}\)\gloss{John 5:24: whoever hears and believes has eternal life and does not come into judgement}
  \breakrow{(11)} \(\pvar{TestimonyArguments.CanonicalWitness}{11}\)\gloss{John 6:28–29: the work of God is that you believe in him whom he has sent}
  \breakrow{(12)} \(\pvar{TestimonyArguments.CanonicalWitness}{12}\)\gloss{Acts 10:43: everyone who believes in him receives forgiveness of sins}
  \breakrow{(13)} \(\pvar{TestimonyArguments.CanonicalWitness}{13}\)\gloss{Acts 15:10–11: no yoke on the disciples; we are saved through the grace of the Lord Jesus}
  \breakrow{(14)} \(\pvar{TestimonyArguments.CanonicalWitness}{14}\)\gloss{Luke 7:50: your faith has saved you}
  \breakrow{(15)} \(\pvar{TestimonyArguments.CanonicalWitness}{15}\)\gloss{Luke 18:9–14: the tax collector, not the Pharisee who listed his works, went home justified}
  \breakrow{(16)} \((\pvar{TestimonyArguments.CanonicalWitness}{7} \land \pvar{TestimonyArguments.CanonicalWitness}{8}) \rightarrow \pvar{TestimonyArguments.CanonicalWitness}{21}\)\gloss{\emph{would give:} Faith is necessary: no one is saved without it}
  \breakrow{(17)} \((\pvar{TestimonyArguments.CanonicalWitness}{14} \land \pvar{TestimonyArguments.CanonicalWitness}{15}) \rightarrow \pvar{TestimonyArguments.CanonicalWitness}{22} \land \pvar{TestimonyArguments.CanonicalWitness}{23}\)\gloss{\emph{would give:} \(\pvar{TestimonyArguments.CanonicalWitness}{22}\) Faith is sufficient: whoever believes is saved; \(\pvar{TestimonyArguments.CanonicalWitness}{23}\) Works are not the ground of justification}
  \item[(18)] \((\pvar{TestimonyArguments.CanonicalWitness}{21} \land \pvar{TestimonyArguments.CanonicalWitness}{22} \land \pvar{TestimonyArguments.CanonicalWitness}{23}) \rightarrow \pvar{TestimonyArguments.CanonicalWitness}{24}\)
  \item[\(\vdash\)] \(\pvar{TestimonyArguments.CanonicalWitness}{24}\)\gloss{\emph{at issue:} Justification is by faith alone}
  \keepwithnext
\end{derivation}
\emph{Grounded in scripture alone:} Romans 3:28: a person is justified by faith apart from works of the law, Ephesians 2:8–9: saved by grace through faith, not of works, Titus 3:5: saved not by works done in righteousness but by his mercy, Romans 4:4–5: to the one who does not work but believes, faith is counted as righteousness, Galatians 2:16: not justified by works of the law but through faith in Christ, Galatians 3:11: no one is justified before God by the law; the righteous shall live by faith, Hebrews 10:38–39: the righteous one shall live by faith, Hebrews 11:6: without faith it is impossible to please God, John 3:16–18: whoever believes has eternal life; whoever does not believe is condemned already, John 5:24: whoever hears and believes has eternal life and does not come into judgement, John 6:28–29: the work of God is that you believe in him whom he has sent, Acts 10:43: everyone who believes in him receives forgiveness of sins, Acts 15:10–11: no yoke on the disciples; we are saved through the grace of the Lord Jesus, Luke 7:50: your faith has saved you, Luke 18:9–14: the tax collector, not the Pharisee who listed his works, went home justified.
\paragraph*{James without its hinge}
\leavevmode\par\nobreak
\par\medskip\noindent\textbf{\texttt{james\allowbreak{}Without\allowbreak{}Demonstrative}} --- James 2\par\nobreak\noindent
James without the demonstrative reading of “justified”: his texts, and the distinction between justification and sanctification, but not the sense of 2:21–24 on which James and Trent part.

\begin{derivation}
  \item[(1)] \(\pvar{TestimonyArguments.CanonicalWitness}{16}\)\gloss{James 2:14–17: faith by itself, if it does not have works, is dead}
  \item[(2)] \(\pvar{TestimonyArguments.CanonicalWitness}{17}\)\gloss{James 2:19: even the demons believe, and shudder}
  \item[(3)] \(\pvar{TestimonyArguments.CanonicalWitness}{18}\)\gloss{James 2:21–23: Abraham’s faith was completed by his works, and Genesis 15:6 was fulfilled}
  \breakrow{(4)} \(\pvar{TestimonyArguments.CanonicalWitness}{19}\)\gloss{James 2:24: a person is justified by works and not by faith alone}
  \breakrow{(5)} \(\pvar{TestimonyArguments.CanonicalWitness}{27}\)\gloss{James 2:14–26 targets a barren faith — mere assent — not Paul’s doctrine}
  \breakrow{(6)} \(\pvar{TestimonyArguments.CanonicalWitness}{29}\)\gloss{Justification and sanctification are inseparable but distinct}
  \item[(7)] \((\pvar{TestimonyArguments.CanonicalWitness}{16} \land \pvar{TestimonyArguments.CanonicalWitness}{18} \land \pvar{TestimonyArguments.CanonicalWitness}{27} \land \pvar{TestimonyArguments.CanonicalWitness}{28} \land \pvar{TestimonyArguments.CanonicalWitness}{29}) \rightarrow \pvar{TestimonyArguments.CanonicalWitness}{25}\)
  \item[\(\vdash\)] \(\pvar{TestimonyArguments.CanonicalWitness}{25}\)\gloss{\emph{at issue:} Good works are the fruit and evidence of a true and lively faith}
\end{derivation}
\par\medskip\noindent\textbf{\texttt{with\allowbreak{}James\allowbreak{}Undemonstrated}} --- The canonical witness with James, without the demonstrative reading\par\nobreak\noindent
The canonical witness with James, on that reading.

\begin{derivation}
  \item[(1)] \(\pvar{TestimonyArguments.CanonicalWitness}{1}\)\gloss{Romans 3:28: a person is justified by faith apart from works of the law}
  \item[(2)] \(\pvar{TestimonyArguments.CanonicalWitness}{5}\)\gloss{Ephesians 2:8–9: saved by grace through faith, not of works}
  \item[(3)] \(\pvar{TestimonyArguments.CanonicalWitness}{6}\)\gloss{Titus 3:5: saved not by works done in righteousness but by his mercy}
  \breakrow{(4)} \(\pvar{TestimonyArguments.CanonicalWitness}{2}\)\gloss{Romans 4:4–5: to the one who does not work but believes, faith is counted as righteousness}
  \breakrow{(5)} \(\pvar{TestimonyArguments.CanonicalWitness}{3}\)\gloss{Galatians 2:16: not justified by works of the law but through faith in Christ}
  \breakrow{(6)} \(\pvar{TestimonyArguments.CanonicalWitness}{4}\)\gloss{Galatians 3:11: no one is justified before God by the law; the righteous shall live by faith}
  \breakrow{(7)} \(\pvar{TestimonyArguments.CanonicalWitness}{7}\)\gloss{Hebrews 10:38–39: the righteous one shall live by faith}
  \breakrow{(8)} \(\pvar{TestimonyArguments.CanonicalWitness}{8}\)\gloss{Hebrews 11:6: without faith it is impossible to please God}
  \breakrow{(9)} \(\pvar{TestimonyArguments.CanonicalWitness}{9}\)\gloss{John 3:16–18: whoever believes has eternal life; whoever does not believe is condemned already}
  \breakrow{(10)} \(\pvar{TestimonyArguments.CanonicalWitness}{10}\)\gloss{John 5:24: whoever hears and believes has eternal life and does not come into judgement}
  \breakrow{(11)} \(\pvar{TestimonyArguments.CanonicalWitness}{11}\)\gloss{John 6:28–29: the work of God is that you believe in him whom he has sent}
  \breakrow{(12)} \(\pvar{TestimonyArguments.CanonicalWitness}{12}\)\gloss{Acts 10:43: everyone who believes in him receives forgiveness of sins}
  \breakrow{(13)} \(\pvar{TestimonyArguments.CanonicalWitness}{13}\)\gloss{Acts 15:10–11: no yoke on the disciples; we are saved through the grace of the Lord Jesus}
  \breakrow{(14)} \(\pvar{TestimonyArguments.CanonicalWitness}{14}\)\gloss{Luke 7:50: your faith has saved you}
  \breakrow{(15)} \(\pvar{TestimonyArguments.CanonicalWitness}{15}\)\gloss{Luke 18:9–14: the tax collector, not the Pharisee who listed his works, went home justified}
  \breakrow{(16)} \(\pvar{TestimonyArguments.CanonicalWitness}{16}\)\gloss{James 2:14–17: faith by itself, if it does not have works, is dead}
  \breakrow{(17)} \(\pvar{TestimonyArguments.CanonicalWitness}{17}\)\gloss{James 2:19: even the demons believe, and shudder}
  \breakrow{(18)} \(\pvar{TestimonyArguments.CanonicalWitness}{18}\)\gloss{James 2:21–23: Abraham’s faith was completed by his works, and Genesis 15:6 was fulfilled}
  \breakrow{(19)} \(\pvar{TestimonyArguments.CanonicalWitness}{19}\)\gloss{James 2:24: a person is justified by works and not by faith alone}
  \breakrow{(20)} \(\pvar{TestimonyArguments.CanonicalWitness}{27}\)\gloss{James 2:14–26 targets a barren faith — mere assent — not Paul’s doctrine}
  \breakrow{(21)} \(\pvar{TestimonyArguments.CanonicalWitness}{29}\)\gloss{Justification and sanctification are inseparable but distinct}
  \breakrow{(22)} \(\pvar{TestimonyArguments.CanonicalWitness}{20}\)\gloss{Jude 20–21: keep yourselves in the love of God, waiting for the mercy of our Lord Jesus Christ}
  \breakrow{(23)} \((\pvar{TestimonyArguments.CanonicalWitness}{1} \land \pvar{TestimonyArguments.CanonicalWitness}{5} \land \pvar{TestimonyArguments.CanonicalWitness}{6}) \rightarrow \pvar{TestimonyArguments.CanonicalWitness}{23}\)\gloss{\emph{would give:} Works are not the ground of justification}
  \breakrow{(24)} \((\pvar{TestimonyArguments.CanonicalWitness}{2} \land \pvar{TestimonyArguments.CanonicalWitness}{3} \land \pvar{TestimonyArguments.CanonicalWitness}{4}) \rightarrow \pvar{TestimonyArguments.CanonicalWitness}{21} \land \pvar{TestimonyArguments.CanonicalWitness}{22}\)\gloss{\emph{would give:} \(\pvar{TestimonyArguments.CanonicalWitness}{21}\) Faith is necessary: no one is saved without it; \(\pvar{TestimonyArguments.CanonicalWitness}{22}\) Faith is sufficient: whoever believes is saved}
  \breakrow{(25)} \((\pvar{TestimonyArguments.CanonicalWitness}{7} \land \pvar{TestimonyArguments.CanonicalWitness}{8}) \rightarrow \pvar{TestimonyArguments.CanonicalWitness}{21}\)\gloss{\emph{would give:} Faith is necessary: no one is saved without it}
  \breakrow{(26)} \((\pvar{TestimonyArguments.CanonicalWitness}{9} \land \pvar{TestimonyArguments.CanonicalWitness}{10} \land \pvar{TestimonyArguments.CanonicalWitness}{11}) \rightarrow \pvar{TestimonyArguments.CanonicalWitness}{22} \land \pvar{TestimonyArguments.CanonicalWitness}{21}\)\gloss{\emph{would give:} \(\pvar{TestimonyArguments.CanonicalWitness}{22}\) Faith is sufficient: whoever believes is saved; \(\pvar{TestimonyArguments.CanonicalWitness}{21}\) Faith is necessary: no one is saved without it}
  \breakrow{(27)} \((\pvar{TestimonyArguments.CanonicalWitness}{12} \land \pvar{TestimonyArguments.CanonicalWitness}{13}) \rightarrow \pvar{TestimonyArguments.CanonicalWitness}{22} \land \pvar{TestimonyArguments.CanonicalWitness}{23}\)\gloss{\emph{would give:} \(\pvar{TestimonyArguments.CanonicalWitness}{22}\) Faith is sufficient: whoever believes is saved; \(\pvar{TestimonyArguments.CanonicalWitness}{23}\) Works are not the ground of justification}
  \breakrow{(28)} \((\pvar{TestimonyArguments.CanonicalWitness}{14} \land \pvar{TestimonyArguments.CanonicalWitness}{15}) \rightarrow \pvar{TestimonyArguments.CanonicalWitness}{22} \land \pvar{TestimonyArguments.CanonicalWitness}{23}\)\gloss{\emph{would give:} \(\pvar{TestimonyArguments.CanonicalWitness}{22}\) Faith is sufficient: whoever believes is saved; \(\pvar{TestimonyArguments.CanonicalWitness}{23}\) Works are not the ground of justification}
  \breakrow{(29)} \((\pvar{TestimonyArguments.CanonicalWitness}{16} \land \pvar{TestimonyArguments.CanonicalWitness}{18} \land \pvar{TestimonyArguments.CanonicalWitness}{27} \land \pvar{TestimonyArguments.CanonicalWitness}{28} \land \pvar{TestimonyArguments.CanonicalWitness}{29}) \rightarrow \pvar{TestimonyArguments.CanonicalWitness}{25}\)\gloss{\emph{would give:} Good works are the fruit and evidence of a true and lively faith}
  \breakrow{(30)} \((\pvar{TestimonyArguments.CanonicalWitness}{21} \land \pvar{TestimonyArguments.CanonicalWitness}{22} \land \pvar{TestimonyArguments.CanonicalWitness}{23}) \rightarrow \pvar{TestimonyArguments.CanonicalWitness}{24}\)\gloss{\emph{would give:} Justification is by faith alone}
  \item[(31)] \((\pvar{TestimonyArguments.CanonicalWitness}{24} \land \pvar{TestimonyArguments.CanonicalWitness}{25}) \rightarrow \pvar{TestimonyArguments.CanonicalWitness}{26}\)
  \item[\(\vdash\)] \(\pvar{TestimonyArguments.CanonicalWitness}{26}\)\gloss{\emph{at issue:} Justification is by faith alone, by a faith that is never alone}
  \keepwithnext
\end{derivation}
\emph{Grounded in scripture alone:} Romans 3:28: a person is justified by faith apart from works of the law, Ephesians 2:8–9: saved by grace through faith, not of works, Titus 3:5: saved not by works done in righteousness but by his mercy, Romans 4:4–5: to the one who does not work but believes, faith is counted as righteousness, Galatians 2:16: not justified by works of the law but through faith in Christ, Galatians 3:11: no one is justified before God by the law; the righteous shall live by faith, Hebrews 10:38–39: the righteous one shall live by faith, Hebrews 11:6: without faith it is impossible to please God, John 3:16–18: whoever believes has eternal life; whoever does not believe is condemned already, John 5:24: whoever hears and believes has eternal life and does not come into judgement, John 6:28–29: the work of God is that you believe in him whom he has sent, Acts 10:43: everyone who believes in him receives forgiveness of sins, Acts 15:10–11: no yoke on the disciples; we are saved through the grace of the Lord Jesus, Luke 7:50: your faith has saved you, Luke 18:9–14: the tax collector, not the Pharisee who listed his works, went home justified, James 2:14–17: faith by itself, if it does not have works, is dead, James 2:19: even the demons believe, and shudder, James 2:21–23: Abraham’s faith was completed by his works, and Genesis 15:6 was fulfilled, James 2:24: a person is justified by works and not by faith alone, Jude 20–21: keep yourselves in the love of God, waiting for the mercy of our Lord Jesus Christ.
\paragraph*{Paul on faith, on its own}
\leavevmode\par\nobreak
Paul’s message about faith, apart from what he says about works: faith is necessary, and it is sufficient. Stated as its own position, so it can be set beside James and asked whether the two cohere — and so the answer can say what it rests on.

\par\medskip\noindent\textbf{\texttt{paul\allowbreak{}On\allowbreak{}Faith}} --- Paul: faith is necessary and sufficient\par\nobreak\noindent
\textbf{Paul: faith is necessary and sufficient.} Romans 4:4–5, Galatians 2:16 and 3:11, read as Paul’s own claim about faith.

\begin{derivation}
  \item[(1)] \(\pvar{TestimonyArguments.CanonicalWitness}{2}\)\gloss{Romans 4:4–5: to the one who does not work but believes, faith is counted as righteousness}
  \item[(2)] \(\pvar{TestimonyArguments.CanonicalWitness}{3}\)\gloss{Galatians 2:16: not justified by works of the law but through faith in Christ}
  \item[(3)] \(\pvar{TestimonyArguments.CanonicalWitness}{4}\)\gloss{Galatians 3:11: no one is justified before God by the law; the righteous shall live by faith}
  \item[(4)] \((\pvar{TestimonyArguments.CanonicalWitness}{2} \land \pvar{TestimonyArguments.CanonicalWitness}{3} \land \pvar{TestimonyArguments.CanonicalWitness}{4}) \rightarrow \pvar{TestimonyArguments.CanonicalWitness}{21} \land \pvar{TestimonyArguments.CanonicalWitness}{22}\)
  \item[\(\vdash\)] \(\pvar{TestimonyArguments.CanonicalWitness}{21} \land \pvar{TestimonyArguments.CanonicalWitness}{22}\)\gloss{\emph{at issue:} \(\pvar{TestimonyArguments.CanonicalWitness}{21}\) Faith is necessary: no one is saved without it; \(\pvar{TestimonyArguments.CanonicalWitness}{22}\) Faith is sufficient: whoever believes is saved}
  \keepwithnext
\end{derivation}
\emph{Grounded in scripture alone:} Romans 4:4–5: to the one who does not work but believes, faith is counted as righteousness, Galatians 2:16: not justified by works of the law but through faith in Christ, Galatians 3:11: no one is justified before God by the law; the righteous shall live by faith.
\par\medskip\noindent\textbf{\texttt{james\allowbreak{}Read\allowbreak{}Against\allowbreak{}Sufficiency}} --- James read against the sufficiency of faith\par\nobreak\noindent
\textbf{James read against sufficiency}, as Trent reads him: every text of James, with “that faith” in 2:14 and “not by faith alone” in 2:24 taken of faith as such. Then faith does not suffice.

\begin{derivation}
  \item[(1)] \(\pvar{TestimonyArguments.CanonicalWitness}{16}\)\gloss{James 2:14–17: faith by itself, if it does not have works, is dead}
  \item[(2)] \(\pvar{TestimonyArguments.CanonicalWitness}{17}\)\gloss{James 2:19: even the demons believe, and shudder}
  \item[(3)] \(\pvar{TestimonyArguments.CanonicalWitness}{18}\)\gloss{James 2:21–23: Abraham’s faith was completed by his works, and Genesis 15:6 was fulfilled}
  \breakrow{(4)} \(\pvar{TestimonyArguments.CanonicalWitness}{19}\)\gloss{James 2:24: a person is justified by works and not by faith alone}
  \item[(5)] \((\pvar{TestimonyArguments.CanonicalWitness}{16} \land \pvar{TestimonyArguments.CanonicalWitness}{19}) \rightarrow \lnot \pvar{TestimonyArguments.CanonicalWitness}{22}\)
  \item[\(\vdash\)] \(\lnot \pvar{TestimonyArguments.CanonicalWitness}{22}\)\gloss{\emph{at issue, the denial of:} Faith is sufficient: whoever believes is saved}
  \keepwithnext
\end{derivation}
\emph{Grounded in scripture alone:} James 2:14–17: faith by itself, if it does not have works, is dead, James 2:19: even the demons believe, and shudder, James 2:21–23: Abraham’s faith was completed by his works, and Genesis 15:6 was fulfilled, James 2:24: a person is justified by works and not by faith alone.
\subsection*{Appendix: the claims and their sources}
\addcontentsline{toc}{subsection}{Appendix: the claims and their sources}
What the texts say, corpus by corpus; the three parts of \emph{faith alone} — that faith is necessary, that it is sufficient, and that works are not the ground; what James adds, that works are the fruit of saving faith; and the readings on which James and Trent part. See \texttt{Testimony\allowbreak{}Arguments.Canonical\allowbreak{}Witness} for the argument this encodes.

\subsubsection*{A citation for every atom}
Every text is cited to scripture alone. That is the argument’s point, and the generated manifest shows it: every text atom is marked as grounded in scripture alone. The rival grants every one of them — Trent quotes most of them itself — so what the texts say is common ground, and nothing in the argument turns on a citation that only one side accepts.

The readings are where the argument is made, and they are cited as readings: each corpus’s step carries its own rated source (\texttt{Testimony\allowbreak{}Arguments.Canonical\allowbreak{}Witness.Lines}). \texttt{cite} holds the argument’s own citations; \texttt{tridentine\allowbreak{}Cite} overrides the one atom on which Trent parts from it, the sense of “justified” in James 2:24.

\subsubsection*{The claims, numbered}
\begin{longtable}{@{}r >{\raggedright\arraybackslash}p{0.30\textwidth} l >{\raggedright\arraybackslash}p{0.19\textwidth} >{\raggedright\arraybackslash}p{0.21\textwidth}@{}}
\toprule
 & claim & kind & held & source \\
\midrule
\endhead
\pvartarget{TestimonyArguments.CanonicalWitness}{1}$P_{1}$ & Romans 3:28: a person is justified by faith apart from works of the law\textsuperscript{\dag} & textual & Christian, hist.-gramm., consensus & Rom 3:28 \\
\pvartarget{TestimonyArguments.CanonicalWitness}{2}$P_{2}$ & Romans 4:4–5: to the one who does not work but believes, faith is counted as righteousness\textsuperscript{\dag} & textual & Christian, hist.-gramm., consensus & Rom 4:4-5 \\
\pvartarget{TestimonyArguments.CanonicalWitness}{3}$P_{3}$ & Galatians 2:16: not justified by works of the law but through faith in Christ\textsuperscript{\dag} & textual & Christian, hist.-gramm., consensus & Gal 2:16 \\
\pvartarget{TestimonyArguments.CanonicalWitness}{4}$P_{4}$ & Galatians 3:11: no one is justified before God by the law; the righteous shall live by faith\textsuperscript{\dag} & textual & Christian, hist.-gramm., consensus & Gal 3:11 \\
\pvartarget{TestimonyArguments.CanonicalWitness}{5}$P_{5}$ & Ephesians 2:8–9: saved by grace through faith, not of works\textsuperscript{\dag} & textual & Christian, hist.-gramm., consensus & Eph 2:8-9 \\
\pvartarget{TestimonyArguments.CanonicalWitness}{6}$P_{6}$ & Titus 3:5: saved not by works done in righteousness but by his mercy\textsuperscript{\dag} & textual & Christian, hist.-gramm., consensus & Titus 3:5 \\
\pvartarget{TestimonyArguments.CanonicalWitness}{7}$P_{7}$ & Hebrews 10:38–39: the righteous one shall live by faith\textsuperscript{\dag} & textual & Christian, hist.-gramm., consensus & Heb 10:38-39 \\
\pvartarget{TestimonyArguments.CanonicalWitness}{8}$P_{8}$ & Hebrews 11:6: without faith it is impossible to please God\textsuperscript{\dag} & textual & Christian, hist.-gramm., consensus & Heb 11:6 \\
\pvartarget{TestimonyArguments.CanonicalWitness}{9}$P_{9}$ & John 3:16–18: whoever believes has eternal life; whoever does not believe is condemned already\textsuperscript{\dag} & textual & Christian, hist.-gramm., consensus & John 3:16-18 \\
\pvartarget{TestimonyArguments.CanonicalWitness}{10}$P_{10}$ & John 5:24: whoever hears and believes has eternal life and does not come into judgement\textsuperscript{\dag} & textual & Christian, hist.-gramm., consensus & John 5:24 \\
\pvartarget{TestimonyArguments.CanonicalWitness}{11}$P_{11}$ & John 6:28–29: the work of God is that you believe in him whom he has sent\textsuperscript{\dag} & textual & Christian, hist.-gramm., consensus & John 6:28-29 \\
\pvartarget{TestimonyArguments.CanonicalWitness}{12}$P_{12}$ & Acts 10:43: everyone who believes in him receives forgiveness of sins\textsuperscript{\dag} & textual & Christian, hist.-gramm., consensus & Acts 10:43 \\
\pvartarget{TestimonyArguments.CanonicalWitness}{13}$P_{13}$ & Acts 15:10–11: no yoke on the disciples; we are saved through the grace of the Lord Jesus\textsuperscript{\dag} & textual & Christian, hist.-gramm., consensus & Acts 15:10-11 \\
\pvartarget{TestimonyArguments.CanonicalWitness}{14}$P_{14}$ & Luke 7:50: your faith has saved you\textsuperscript{\dag} & textual & Christian, hist.-gramm., consensus & Luke 7:50 \\
\pvartarget{TestimonyArguments.CanonicalWitness}{15}$P_{15}$ & Luke 18:9–14: the tax collector, not the Pharisee who listed his works, went home justified\textsuperscript{\dag} & textual & Christian, hist.-gramm., consensus & Luke 18:9-14 \\
\pvartarget{TestimonyArguments.CanonicalWitness}{16}$P_{16}$ & James 2:14–17: faith by itself, if it does not have works, is dead\textsuperscript{\dag} & textual & Christian, hist.-gramm., consensus & Jas 2:14-17 \\
\pvartarget{TestimonyArguments.CanonicalWitness}{17}$P_{17}$ & James 2:19: even the demons believe, and shudder\textsuperscript{\dag} & textual & Christian, hist.-gramm., consensus & Jas 2:19 \\
\pvartarget{TestimonyArguments.CanonicalWitness}{18}$P_{18}$ & James 2:21–23: Abraham’s faith was completed by his works, and Genesis 15:6 was fulfilled\textsuperscript{\dag} & textual & Christian, hist.-gramm., consensus & Jas 2:21-23 \\
\pvartarget{TestimonyArguments.CanonicalWitness}{19}$P_{19}$ & James 2:24: a person is justified by works and not by faith alone\textsuperscript{\dag} & textual & Christian, hist.-gramm., consensus & Jas 2:24 \\
\pvartarget{TestimonyArguments.CanonicalWitness}{20}$P_{20}$ & Jude 20–21: keep yourselves in the love of God, waiting for the mercy of our Lord Jesus Christ\textsuperscript{\dag} & textual & Christian, hist.-gramm., consensus & Jude 1:20-21 \\
\pvartarget{TestimonyArguments.CanonicalWitness}{21}$P_{21}$ & Faith is necessary: no one is saved without it & theological & Christian, hist.-gramm., consensus & Heb 11:6; John 3:16-18; \cite[Trent, Session VI (1547), Decree on Justification, ch. 8]{tanner-decrees-1990} \\
\pvartarget{TestimonyArguments.CanonicalWitness}{22}$P_{22}$ & Faith is sufficient: whoever believes is saved & theological & Reformed Protestant, disputed & Rom 4:4-5; John 5:24; \cite{schreiner-faith-alone-2015} \\
\pvartarget{TestimonyArguments.CanonicalWitness}{23}$P_{23}$ & Works are not the ground of justification & theological & Reformed Protestant, disputed & Rom 3:28; Eph 2:8-9; \cite{schreiner-faith-alone-2015} \\
\pvartarget{TestimonyArguments.CanonicalWitness}{24}$P_{24}$ & Justification is by faith alone & theological & Reformed Protestant, disputed & \cite[XI.2]{westminster-confession-1647}; \cite[III.xi.19]{calvin-institutes-1960}; \cite[8.6]{hilary-commentary-matthew-2013}; \cite[Homily 20, Of Humility]{basil-ascetical-works-1999} \\
\pvartarget{TestimonyArguments.CanonicalWitness}{25}$P_{25}$ & Good works are the fruit and evidence of a true and lively faith & theological & Reformed Protestant, well supported & \cite[XVI.2]{westminster-confession-1647} \\
\pvartarget{TestimonyArguments.CanonicalWitness}{26}$P_{26}$ & Justification is by faith alone, by a faith that is never alone & theological & Reformed Protestant, disputed & \cite[XI.2]{westminster-confession-1647} \\
\pvartarget{TestimonyArguments.CanonicalWitness}{27}$P_{27}$ & James 2:14–26 targets a barren faith — mere assent — not Paul’s doctrine & interpretive & Reformed Protestant, well supported & \cite[ad loc. Jas 2:14]{moo-james-2000}; \cite[ad loc. Jas 2:24]{johnson-james-1995}; Jas 2:19 \\
\pvartarget{TestimonyArguments.CanonicalWitness}{28}$P_{28}$ & ‘Justified’ in James 2:21–24 means shown to be righteous (denied by Trent, which reads 2:24 as an increase of justification by works done in grace) & linguistic & Roman Catholic, disputed & \cite[Trent, Session VI (1547), Decree on Justification, ch. 10]{tanner-decrees-1990}; Jas 2:24 \\
\pvartarget{TestimonyArguments.CanonicalWitness}{29}$P_{29}$ & Justification and sanctification are inseparable but distinct & theological & Reformed Protestant, disputed & \cite[Q. 77]{westminster-larger-catechism-1647}; 1 Cor 6:11; 1 Cor 1:30; Rom 4:6; Rom 4:8; Ezek 36:27; Heb 9:13-14; Rom 3:24-25; Rom 6:6; Rom 6:14; \cite[III.xi.6]{calvin-institutes-1960}; \cite[Epistle to Diognetus 9]{roberts-ante-nicene-fathers-1-1885} \\
\bottomrule
\end{longtable}

\noindent\footnotesize \textsuperscript{\dag} grounded in scripture alone: the reading is assumed, not argued for.
\normalsize
\subsubsection*{How this is checked}
\par\medskip\noindent\textbf{\texttt{text\allowbreak{}Only}}\par\nobreak\noindent
A claim about what a text says, cited to the text alone. Every party to the argument accepts the text as scripture, so the claim is \texttt{consensus}.

{\small\begin{verbatim}
def textOnly : List ScriptureCitation → Source :=
  fun refs =>
    { primary := Reference.scripture refs,
      tradition := Tradition.christianHistoricalGrammatical,
      confidence := Confidence.consensus }
\end{verbatim}}
\par\medskip\noindent\textbf{\texttt{verse}}\par\nobreak\noindent
A verse, as a citation.

{\small\begin{verbatim}
def verse : Book → ℕ → ℕ → ScriptureCitation :=
  fun b c v =>
    { ref := PassageRange.verse { book := b, chapter := c, verse := v } }
\end{verbatim}}
\par\medskip\noindent\textbf{\texttt{verses}}\par\nobreak\noindent
A range of verses in one chapter, as a citation.

{\small\begin{verbatim}
def verses : Book → ℕ → ℕ → ℕ → ScriptureCitation :=
  fun b c v w =>
    {
      ref :=
        PassageRange.range
          { book := b, startChapter := c, startVerse := v, endChapter := c,
            endVerse := w } }
\end{verbatim}}
\par\medskip\noindent\textbf{\texttt{cite}}\par\nobreak\noindent
The argument’s citation of every atom.

{\small\begin{verbatim}
def cite : Claim → AtomMeta
\end{verbatim}}
\par\medskip\noindent\textbf{\texttt{tridentine\allowbreak{}Cite}}\par\nobreak\noindent
Trent’s citations: the argument’s, except the one atom Trent denies, the demonstrative sense of “justified” in James 2:24. Trent reads it as the increase of a justification already received, quoting the verse. The label states the claim Trent denies, so that “not” before it reads correctly.

{\small\begin{verbatim}
def tridentineCite : Claim → AtomMeta
\end{verbatim}}

\bibliographystyle{plain}
\bibliography{references}

\clearpage
\section*{How this edition was made}
\addcontentsline{toc}{section}{How this edition was made}
This document was printed by a program, not set by hand. Its text comes from the source files of the Testimony library, where each argument is written out twice: in prose, for a reader, and in Lean, a language for stating proofs that a computer can check. The prose is printed in the order it is written, and under each part come the premises, lines of reasoning and results it introduces, as the source states them. Before anything is printed, the computer checks every one of those results: where this document says a conclusion follows from its premises, that it does, and where it says a conclusion does not follow, that the counterexample it gives really is one. It does not check that any premise is true.

\ifeditionrecorded This edition, version \editionversion, was made from the library as it stood on \editiondate, at the commit named in step 3 below. A commit is one saved state of every file in the repository; that string of letters and digits names it, and what it names never changes. So anyone can make this edition again.
\else This copy was typeset without a record of the commit it was made from, so the steps below write \texttt{COMMIT} where an edition names its own. A commit is one saved state of every file in the repository, named by a string of letters and digits that never changes what it names.
\fi

You need a computer running Linux or macOS (on Windows, the Windows Subsystem for Linux), with git, Python~3, an internet connection, and about 10\,GB of free disk space.
\begin{enumerate}
\item Install elan, which installs the version of Lean the library asks for:
\begin{alltt}\small
curl https://elan.lean-lang.org/elan-init.sh -sSf | sh
\end{alltt}
\item Install Tectonic, which sets the pages, following \url{https://tectonic-typesetting.github.io/book/latest/installation/}.
\item Fetch the library, and turn it back to this edition's commit:
\begin{alltt}\small
git clone https://github.com/testimonyproject/testimony.git
cd testimony
git checkout \editioncommit
\end{alltt}
\item Download the mathematics library Lean builds on, already compiled, so it need not be compiled on your computer:
\begin{alltt}\small
lake exe cache get
\end{alltt}
\item Check every argument. This is the longest step, and if any proof failed it would stop here with an error:
\begin{alltt}\small
cd arguments
lake build
\end{alltt}
\item Write the document from the source:
\begin{alltt}\small
lake exe argtex
cd ..
\end{alltt}
\item Record the edition, and work out from the repository's history who wrote each argument:
\begin{alltt}\small
python3 scripts/credits.py docs/latex docs/latex/arguments
python3 scripts/edition.py docs/latex docs/latex/arguments
\end{alltt}
\item Set the pages:
\begin{alltt}\small
cd docs/latex/arguments
tectonic canonical-witness.tex
\end{alltt}
\end{enumerate}
That makes \texttt{canonical-witness.pdf}. Its words, numbers and references are this edition's, and so should its pages be: the fonts it is set in come with Tectonic, not from your computer.
\end{document}
